\documentclass[11pt,letterpaper]{article}
\usepackage[T1]{fontenc}
\usepackage[utf8]{inputenc}
\usepackage{lmodern}
\usepackage[margin=1in]{geometry}
\usepackage{amsmath,amssymb,amsthm}
\usepackage{microtype}
\usepackage{booktabs,longtable,array}
\usepackage{listings}
\usepackage{xurl}
\usepackage[hidelinks]{hyperref}

\hypersetup{pdftitle={Moment and Prime-Trace Criteria for the Riemann Hypothesis, with Two Spectral Exclusions},
 pdfauthor={Henry Arellano-Pe\~na},
 pdfsubject={Moment and trace criteria, spectral exclusions, and a TCGS-SEQUENTION realization problem}}
\setlength{\emergencystretch}{2em}
\setlength{\parskip}{0.25em}
\numberwithin{equation}{section}
\newtheorem{theorem}{Theorem}[section]
\newtheorem{proposition}[theorem]{Proposition}
\newtheorem{lemma}[theorem]{Lemma}
\newtheorem{corollary}[theorem]{Corollary}
\theoremstyle{definition}
\newtheorem{definition}[theorem]{Definition}
\newtheorem{assumption}[theorem]{Assumption}
\theoremstyle{remark}
\newtheorem{remark}[theorem]{Remark}
\newcommand{\R}{\mathbb R}
\newcommand{\C}{\mathbb C}
\newcommand{\N}{\mathbb N}
\newcommand{\CS}{\mathcal C}
\newcommand{\HS}{\mathcal H}
\newcommand{\ES}{\mathcal E_{\mathcal S}}
\newcommand{\src}{\boldsymbol\Xi}
\newcommand{\Id}{\mathrm I}
\newcommand{\dd}{\,\mathrm d}
\newcommand{\e}{\mathrm e}
\DeclareMathOperator{\Tr}{Tr}
\DeclareMathOperator{\Dom}{Dom}
\DeclareMathOperator{\Dim}{Dim}
\DeclareMathOperator{\Ran}{Ran}
\DeclareMathOperator{\rank}{rank}
\DeclareMathOperator{\supp}{supp}
\DeclareMathOperator{\Info}{Info}
\DeclareMathOperator{\Read}{Read}
\DeclareMathOperator{\ev}{ev}
\DeclareMathOperator{\Orb}{Orb}
\DeclareMathOperator{\Li}{Li}
\DeclareMathOperator{\arcosh}{arcosh}
\newcolumntype{P}[1]{>{\raggedright\arraybackslash}p{#1}}
\lstset{basicstyle=\ttfamily\footnotesize,breaklines=true,
 columns=fullflexible,keepspaces=true,showstringspaces=false,
 frame=single,rulecolor=\normalcolor,numbers=left,numberstyle=\tiny,
 xleftmargin=1.5em,framexleftmargin=1.2em}

\title{\textbf{Moment and Prime-Trace Criteria for the Riemann Hypothesis, with Two Spectral Exclusions}}
\author{Henry Arellano-Pe\~na\\
\small Nuevo Estandar Biotropical NEBIOT S.A.S., Bogot\'a, Colombia\\
\small \href{mailto:harellano@unal.edu.co}{harellano@unal.edu.co}}
\date{6 September 2026}

\begin{document}
\maketitle
\vspace{-1.5em}
\begin{abstract}
We develop moment and prime-trace formulations of the Riemann hypothesis and prove two classes of spectral exclusions. The analytic results use explicitly stated entire-function, moment, and operator hypotheses; a separate TCGS--SEQUENTION application formulates the source-to-readout realization problem. Writing $F(w)=\xi(\tfrac12+i\sqrt w)/\xi(\tfrac12)$ by its entire power series, we establish a self-contained chain between the Riemann hypothesis, positivity of every Hankel matrix generated by $-F'/F$, and a positive trace-class Fredholm determinant. The moment formulation is situated in the classical Grommer--Hamburger real-zero theory, with its reciprocal-square indexing made explicit. Its connections to typed constitutive readout and selector covariance are conditional realization statements. A Gram-transfer theorem identifies the exact arithmetic identity that a linearized Extrinsic Constitutive Law must satisfy, and a finite-rank obstruction excludes any single finite protected response sector from representing the full moment tower. An equivalent prime-side resolvent identity is derived in the absolutely convergent half-plane, avoiding any use of the Euler product inside the critical strip. The required counting and heat-trace asymptotics exclude the full spectrum of ordinary fixed-order elliptic operators on smooth compact manifolds. We then examine an explicit exponential-corridor Schr\"odinger operator: it is positive and self-adjoint, has a trace-class inverse, and reproduces both leading terms of the zeta-zero counting law. Nevertheless, a large-order Bessel calculation proves that its normalized determinant differs from $F$ by a ratio asymptotic to a positive constant times $s^{9/4}$ on the real $s$-axis. Allowing both the Morse potential scale and the Dirichlet endpoint to vary does not repair the determinant: the second counting coefficient fixes their spectral combination to $4\pi$, the power prefactor forces $\kappa=9/4$, and a nonzero inverse-order coefficient still excludes equality. An explicit complex-order estimate establishes the Whittaker genus-zero determinant normalization. The leading counting term also fixes the kinetic--rate invariant of the full two-exponential Morse family. Constant spectral shifts introduce an unmatched logarithmic inverse-order term. Fixed Robin data change the power match to $\kappa=5/4$ and can cancel the first inverse-order discrepancy, but a nonzero second-order coefficient excludes the remaining candidate without a numerical normalization premise. Finite numerical moment tests and a counterexample family clarify why symmetry, kernel positivity, and truncated positivity cannot close the argument. Compatible cochains and a shared boundary law motivate a positive Schur-complement response. We distinguish its static determinant from the full spectral pencil, prove finite-to-infinite moment transfer and a trace-tightness estimate, and derive explicit prime-tail and resolvent-trace error bounds for candidate and bounded-family rejection. The Riemann hypothesis is not proved; the remaining step is precisely the construction of a source-derived positive representation satisfying the full arithmetic Gram or prime-trace identity.
\end{abstract}

\begin{samepage}
{\small \noindent\textbf{Keywords:} Riemann hypothesis; TCGS--SEQUENTION; 4D-Counterduction; Counterspace; Extrinsic Constitutive Law; source--shadow geometry; Hamburger moment problem; Fredholm determinant; spectral geometry; prime numbers; selector invariance; conservative cochains; boundary response; Schur complement; trace-norm convergence.\par}
\end{samepage}

\section{Introduction}
\label{sec:introduction}
The Riemann hypothesis concerns the zeros of a particular analytic function, not an unspecified resemblance between order and irregularity. Its arithmetic importance comes from the fact that the Riemann zeta function simultaneously has a Dirichlet-series representation over the positive integers and an Euler-product representation over the primes. Riemann's analytic continuation and explicit-formula programme connect these descriptions to the location of complex zeros \cite{riemann,titchmarsh}. The hypothesis remains listed as unresolved by the Clay Mathematics Institute \cite{clay}. That status fixes the question to be addressed; it does not determine which mathematical constructions should be attempted.

TCGS--SEQUENTION supplies a different starting language from conventional arithmetic: a complete non-temporal four-dimensional source, a three-dimensional operational Shadow, a conserved singular structure, and an Extrinsic Constitutive Law (ECL) governing admissible readout. Its geometric and operator formulations distinguish a source configuration from its informational organization, the complete readability profile from one selected value, and an observable from the mathematical map used to represent it \cite{tcgsfoundation,tcgsgeometry,tcgsoperator}. An application to number theory must retain these distinctions. In particular, neither an integer nor a complex zero is identified with a source singularity merely because the framework uses singular structure elsewhere.

The central research question is therefore constructive: can a source-grounded geometric representation produce a positive operator whose determinant or trace is \emph{exactly} the relevant arithmetic function? Self-adjointness alone is insufficient. There are innumerable self-adjoint operators with spectra unrelated to prime numbers. Conversely, an operator defined retrospectively from zeros already assumed to be real cannot serve as an independent proof of their reality. The arithmetic identification and the positivity mechanism must be established without one being concealed inside the definition of the other. This is the central difficulty of spectral approaches, including the Hilbert--P\'olya perspective and its semiclassical developments \cite{berrykeating}.

We pursue two complementary constructions. The first extracts a moment geometry directly from the completed zeta function. Its positive realization would settle the zero-location question. The second starts with an explicit positive geometric operator and asks whether its spectral determinant equals the arithmetic target. These directions meet at an exact trace identity. The first construction supplies arithmetic specificity but leaves universal positivity open; the second supplies positivity and the correct leading spectral density but fails arithmetic specificity. Their comparison identifies a concrete missing bridge rather than an unspecified demand for further formalization.

The moment approach belongs to an established mathematical landscape. Li's criterion expresses the hypothesis through positivity of a different sequence \cite{li}; Cardon and Roberts relate a Riemann-hypothesis condition, with simplicity, to orthogonal polynomials \cite{cardon}; Voros develops zeta functions over zeta zeros \cite{voros}. The real-zero/Hankel method originates in Grommer's criterion \cite{grommer}, whose inverse-power indexing is explicitly restated by Baricz and \v{S}tampach \cite[Theorem~10]{bariczstampach}. Levin provides the general entire-function setting \cite{levin}; related Laguerre and Tur\'an inequalities concern different positivity conditions \cite{cravencsordas}. Section~\ref{subsec:grommernormalization} gives the precise even/odd block relation to the moments used here rather than identifying the differently indexed criteria by name alone. The moment equivalence below is not claimed as a new discovery. It is proved in the normalization needed for the present source-to-readout construction, with the dependence of every subsequent result made explicit.

The proved contribution is the resulting mathematical architecture and its executed tests. These proofs use classical analysis and explicit linear-algebraic hypotheses, not the TCGS axioms as additional mathematical premises. The framework supplies the realization question and its source/readout restrictions; Assumption~\ref{ass:gram} remains an unconstructed arithmetic identification. We obtain an ECL-to-arithmetic Gram-transfer theorem; a finite-rank obstruction; a prime-side trace identity sufficient and necessary for a positive determinant realization; a compact-elliptic spectral exclusion; and an explicit determinant mismatch for a calibrated exponential model. The exponential construction is a rescaled member of the classical Bessel--Morse family, not a newly discovered solvable potential \cite{lagarias}. The substantive calculation is the comparison with the fully normalized Riemann function, including the term that prevents equality despite agreement of the leading counting law. The expanded comparison includes all fixed real separated endpoint conditions and constant spectral shifts; a controlled second-order coefficient closes the Robin case after its first-order matching succeeds. The leading counting term also fixes the product of the exponential rate and the square root of the kinetic coefficient, extending the exclusion to the full constant-coefficient Morse family with rate ratio $2:1$.

This distinction also determines the role of the framework. Non-temporal Counterspace is the governing starting point, not a hypothesis tested by the numerical signs in this paper. The object exposed to mathematical failure is the \emph{particular arithmetic realization}. A failed potential or an unproved Gram identity cannot be repaired by restating source completeness. Equally, rejecting one potential does not establish a no-go theorem for every realization compatible with the framework. The paper carries out the local construction, proves what it proves, and specifies the exact condition it does not obtain.

The conservative source-to-readout and reciprocal fractal-foam constructions provide a further route through compatible cochains, shared singular-boundary data, and interior elimination \cite{tcgsconservative,tcgsfoam}. We formulate its positive interface response without identifying it with the full bulk determinant, and connect finite refinements to arithmetic positivity through an all-degree limit theorem. Quantitative trace and prime-tail estimates then give certified rejection conditions. These additions execute the algebraic reduction and limit implications, while leaving the source-derived arithmetic identification explicit.

\section{Source geometry and the arithmetic realization problem}
\label{sec:framework}
\subsection{The governing architecture}
The four governing requirements are stated in their source--shadow order. Whole Content requires a complete non-temporal four-dimensional Counterspace, represented geometrically by $(\CS^4,G,\Psi)$. Identity of Source supplies a distinguished origin $p_0$ and a conserved singular set
\begin{equation}
 \mathcal S=\Orb(p_0).
 \label{eq:singular}
\end{equation}
Shadow realization supplies an admissible immersion $X:\Sigma^3\hookrightarrow\CS^4$ and pullback data $(g,\psi)=X^*(G,\Psi)$. Parsimony requires one source-grounded constitutive/readability architecture rather than independently adjustable source laws for each application \cite{tcgsfoundation}.

The foliation, its label, and operational clock time have distinct roles. An admissible foliation $\mathcal F$ is geometric readout structure. Its ordering label $\lambda_g$ is defined only up to admissible monotone reparameterization. Operational time $\tau_{\mathrm{op}}$ is a functional of Shadow clocks, correlations, and records. Thus
\begin{equation}
 \mathcal F\not\equiv\lambda_g\not\equiv\tau_{\mathrm{op}},
 \qquad t_{\mathrm{ont}}\notin\Dim(\CS^4).
 \label{eq:time}
\end{equation}
Here $\Dim(\CS^4)$ denotes the catalogue of source dimensions rather than an additional mathematical operator. The notation prevents an analytical parameter introduced later from being mistaken for an ontic temporal direction \cite{tcgsfoundation}.

\subsection{Law, profile, readout, and observable}
Let $\src\in\mathcal X_{\mathrm{adm}}$ denote an admissible source configuration. Boldface distinguishes this source datum from the conventional Riemann function $\Xi(z)$. The typed informational map is
\begin{equation}
 \Info:\mathcal X_{\mathrm{adm}}\longrightarrow\mathcal I_{\mathrm{src}},
 \qquad M_{\src}=\Info(\src).
\end{equation}
For selectors $\sigma=(X,[s],\mathcal B)$, let $\mathcal Y_\sigma$ be the corresponding readout spaces and $\mathcal P_{\mathrm{read}}=\prod_\sigma\mathcal Y_\sigma$. The ECL and its evaluations are
\begin{align}
 \ES&=\Read_{\mathcal S,G}\circ\Info,
       &\ES&:\mathcal X_{\mathrm{adm}}\to\mathcal P_{\mathrm{read}},\label{eq:ecl}\\
 \mathcal E_{\mathcal S,\src}&=\ES(\src),
       &\mathcal E^\sigma_{\mathcal S,\src}&=\ev_\sigma\mathcal E_{\mathcal S,\src}.
\end{align}
An independently specified physical-readout map $\pi^{\mathrm{phys}}_\sigma$ then gives
\begin{equation}
 \mathcal O^\sigma_{\mathcal S}
 =\pi^{\mathrm{phys}}_\sigma\circ\ev_\sigma\circ\Read_{\mathcal S,G}\circ\Info.
 \label{eq:counterduction}
\end{equation}
This composition is the relevant 4D-Counterduction relation. It is not Counterspace itself, a second source manifold, or an independently evolving source process \cite{tcgsfoundation,tcgsoperator}.

In particular, the map $\ES$, its profile value, and a scalar function subsequently used in an arithmetic calculation are not interchangeable. The scalar function $F$ introduced below is an analytic arithmetic target. No equation $F=\ES$ is asserted. Such an equality would erase both the source dependence of the ECL and the domain and codomain of the arithmetic function. An arithmetic application must instead specify a representation connecting an admissible source reduction to $F$, while retaining the distinction between the mathematical representation and the source it represents \cite{tcgsfoundation}.

\subsection{The legitimate positive operator}
The relevant operator construction is based on a physical-readout derivative, not an untyped Hessian of a profile-valued law. At a reference source configuration $\src_0$, suppose
\begin{equation}
 L_{\sigma,0}=D\mathcal O^\sigma_{\mathcal S}[\src_0]
\end{equation}
is densely defined and closable between specified Hilbert completions. Writing $L_\sigma$ for its closure, the natural operator
\begin{equation}
 O_\sigma=L_\sigma^*L_\sigma,\qquad
 \Dom(O_\sigma)=\{v\in\Dom(L_\sigma):L_\sigma v\in\Dom(L_\sigma^*)\}
 \label{eq:observability}
\end{equation}
is positive and self-adjoint. Its quadratic form is
\begin{equation}
 \langle v,O_\sigma v\rangle=\|L_\sigma v\|^2.
 \label{eq:readoutnorm}
\end{equation}
The mathematical statement is standard closed-operator theory \cite{reed}; its use here follows the explicitly typed observability construction in the framework \cite{tcgsoperator}.

The interpretation matters. Equation~\eqref{eq:readoutnorm} measures squared readout sensitivity in the declared norm. It does not automatically measure energy, nor does it imply that $O_\sigma$ has a discrete spectrum, a trace-class inverse, or the zeta zeros as eigenvalues. Banach charts, Hilbert metrics, derivative domains, closability, and selector comparison maps are additional analytic realization data. They are not inferred from the number of source dimensions \cite{tcgsoperator}.

\subsection{What arithmetic must add}
The formal inputs $(\CS,G,\Psi,p_0,\mathcal S)$ do not yet specify an action of integer multiplication, an Euler product, a prime-weighted trace, or a polynomial Gram map. Introducing such a map is legitimate as a proposed mathematical realization, but its construction must be displayed. A finite-dimensional protected response sector can be useful for a restricted physical application without encoding an infinite arithmetic spectrum. Section~\ref{sec:gram} makes this difference a rank obstruction rather than a verbal caution.

Nor is a one-dimensional analytical coordinate prohibited by a four-dimensional source ontology. A coordinate on an auxiliary function space, a reduced corridor, or a Mellin transform is not a claim that the physical source has become one-dimensional. The distinction follows the framework's separation of source dimension, representation dimension, and operational carrier \cite{tcgsfoundation,tcgsgeometry}. It permits the half-line model below while preventing its formal success from identifying that half-line with Counterspace.

\subsection{Conservative source cartography and reciprocal interfaces}
\label{subsec:sourcecompatible}
Two geometric constructions give additional content to the realization problem: the conservative discrete source-to-readout architecture and the reciprocal fractal-foam companion \cite{tcgsconservative,tcgsfoam}. They supply distinct ingredients. The former specifies orientation, dual volumes, shared-flux cancellation, compatible cochains, boundary completion, and selector transport. The latter supplies a candidate phase field $\phi$, reciprocal closed two-form, nonlinear potential $U$, and source-determined interface and corridor weights. Neither construction currently supplies an arithmetic insertion map or a prime trace. Their use below is restricted to those structures they actually define.

The conservative paper distinguishes the ontic ECL from its scientific representation by a hat,
\begin{equation}
 \widehat{\mathcal E}_{\mathcal S,G}
 =\Read_{\mathcal S,G}\circ\Info.
 \label{eq:hattedcartography}
\end{equation}
The symbol $\ES$ in Equation~\eqref{eq:ecl} is used at this same formal-map level throughout the present analysis. No finite operator, equation, or mesh is identified exhaustively with the ontic law. This clarification preserves the map/profile/readout types while adopting the explicit distinction in the two geometric papers \cite{tcgsconservative,tcgsfoam}.

Let $r$ index a declared refinement sequence on regularized source domains
\begin{equation}
 \Omega_{\varrho_r}=\CS\setminus N_{\varrho_r}(\mathcal S),
 \qquad h_r\downarrow0,\quad\varrho_r\downarrow0.
 \label{eq:regularizedsource}
\end{equation}
The sequence may additionally resolve diffuse interfaces or expand a noncompact domain; every such limit and its order must be specified. At fixed $r$, conservative assembly requires a shared flux with $\widehat Q_{kj}=-\widehat Q_{jk}$. A reciprocal-form realization additionally requires
\begin{equation}
 d_{q+1,r}d_{q,r}=0,
 \qquad d_rX_{\sigma,r}^{*}=X_{\sigma,r}^{*}d_r.
 \label{eq:compatiblecochains}
\end{equation}
The nodal volumes and directed hyperareas do not supply the full cochain complex or its Hodge matrices. They cannot be used as though the reciprocal two-form were an ordinary scalar nodal variable \cite{tcgsconservative,tcgsfoam}. Stable approximation of the associated Hilbert complex requires suitable subcomplexes and bounded commuting projections in addition to the algebraic identities \cite{arnoldfeec}.

Gate T of the geometric constructions fixes the source action, singular orbit, complement topology, periods, and induced carrier classes before a harmonic sector is admitted. For example, the isolated-point complement in a four-ball has trivial degree-two cohomology. Gate S then requires all singular-boundary channels to descend from one declared functional $\mathcal B_{\mathcal S}$, with coefficients and excision scaling fixed independently of the target. Its prescribed variation jointly generates the phase flux, nonlinear-potential flux, reciprocal current, and traction:
\begin{equation}
 \delta\mathcal B_{\mathcal S}
 =\int_{\partial N_{\varrho_r}(\mathcal S)}
 \left[-Q^\phi_{\mathcal S}\delta\phi-Q^U_{\mathcal S}\delta U
 +\langle J_{\mathcal S},\delta A_{\parallel}\rangle
 -\langle t_{\mathcal S},\delta X\rangle\right]\dd A_G.
 \label{eq:sharedboundaryvariation}
\end{equation}
The geometric papers state this requirement but do not construct the required shared functional. It therefore enters the response construction as explicit missing model data, not as a completed source derivation \cite{tcgsconservative,tcgsfoam}.

The companion's corridor metric is $\widetilde G=w_{\mathrm{cor}}^{-2}G$, where $w_{\mathrm{cor}}$ denotes its source-determined readability weight, distinguished here from the complex arithmetic variable $w$. Its multifractal weight spectrum is not an eigenvalue counting theorem. Neither a corridor nor a harmonic period may be assigned a prime label and then cited as a derivation of that label. The reported median-dual and manufactured-curvature checks are implementation evidence within their stated tests; they do not establish spectral convergence, an arithmetic determinant, or the unresolved boundary functional \cite{tcgsconservative,tcgsfoam}.

\section{The arithmetic target and its entire normalization}
\label{sec:normalization}
For $\Re s>1$, define
\begin{equation}
 \zeta(s)=\sum_{n=1}^{\infty}n^{-s}
 =\prod_{p\ \mathrm{prime}}(1-p^{-s})^{-1}.
 \label{eq:euler}
\end{equation}
The domain in Equation~\eqref{eq:euler} is essential: its absolutely convergent Euler product is not used on the critical line. Analytic continuation gives the meromorphic zeta function, and the completion
\begin{equation}
 \xi(s)=\frac12s(s-1)\pi^{-s/2}\Gamma(s/2)\zeta(s)
 \label{eq:xi}
\end{equation}
is entire of order one. It satisfies $\xi(s)=\xi(1-s)$ and $\xi(\overline s)=\overline{\xi(s)}$. Its zeros are precisely the nontrivial zeta zeros, with multiplicity \cite{riemann,titchmarsh,dlmf}.

Set
\begin{equation}
 \Xi(z)=\xi\!\left(\frac12+iz\right),\qquad
 F(w)=\frac{\Xi(\sqrt w)}{\Xi(0)}.
 \label{eq:normal}
\end{equation}
The second expression is notation for an entire power series, not a branch choice. Since $\Xi$ is even,
\begin{equation}
 \Xi(z)=\sum_{j=0}^{\infty}c_jz^{2j},\qquad
 F(w)=\frac1{c_0}\sum_{j=0}^{\infty}c_jw^j.
\end{equation}
Consequently $F$ is a real entire function of order $1/2$ and $F(0)=1$. If $\rho=\beta+i\gamma$ is a zero with $\gamma>0$, its corresponding zero of $F$ is
\begin{equation}
 w_\rho=-(\rho-\tfrac12)^2
       =\bigl(\gamma+i(\tfrac12-\beta)\bigr)^2.
 \label{eq:zerochange}
\end{equation}
Exactly one representative of the pair $\rho,1-\rho$ is counted by taking $\gamma>0$. Off-line zeros give conjugate nonreal values of $w_\rho$. On-line zeros give $w_\rho=\gamma^2>0$.

\begin{lemma}[Genus-zero normalization]
\label{lem:product}
Enumerate the zeros of $F$ as $(w_j)$ with multiplicity. Then
\begin{equation}
 \sum_j|w_j|^{-1}<\infty,\qquad
 F(w)=\prod_j\left(1-\frac w{w_j}\right).
 \label{eq:product}
\end{equation}
Moreover $F(-x)>0$ for $x\geq0$.
\end{lemma}
\begin{proof}
The standard zero-counting estimate $N(T)=O(T\log T)$, where $N(T)$ counts nontrivial zeros with $0<\Im\rho\leq T$, implies $\sum_{\Im\rho>0}|\rho-1/2|^{-2}<\infty$ \cite{titchmarsh}. This is the first statement. Since $F$ has order $1/2<1$, its Hadamard factorization contains no nonconstant exponential factor. The normalization $F(0)=1$ gives Equation~\eqref{eq:product}. Finally, Section~\ref{sec:theta} derives the positive-kernel formula
\[
 \xi(\tfrac12+r)=\int_0^\infty\Phi(u)\cosh(ru)\dd u>0
 \quad(r\in\R).
\]
Taking $r=\sqrt x$ and using the functional equation proves $F(-x)>0$.
\end{proof}

\begin{corollary}
\label{cor:positivezeros}
The Riemann hypothesis is equivalent to every zero of $F$ being positive real.
\end{corollary}
\begin{proof}
Under the hypothesis Equation~\eqref{eq:zerochange} gives $w_\rho=\gamma^2$. Conversely, if $w_\rho>0$, the number $z=\gamma+i(1/2-\beta)$ has positive real part and real positive square. Its imaginary part is therefore zero, giving $\beta=1/2$.
\end{proof}

This normalization removes an avoidable ambiguity in spectral programmes. A self-adjoint operator with eigenvalues $\gamma$ is usually discussed on both sides of zero. A positive trace-class operator will instead have eigenvalues $\gamma^{-2}$. These are different spectral variables, with different counting laws and operator classes. Keeping the square and inverse explicit prevents an asymptotic calculation in one variable from being mistaken for a result in another.

\section{Why the critical half is natural but not sufficient}
\label{sec:half}
\subsection{Reflection does not force a zero onto its fixed set}
The holomorphic map $s\mapsto1-s$ has the single fixed point $s=1/2$. The antiholomorphic reflection $s\mapsto1-\overline s$ has the critical line as its fixed set. The functional equation and complex conjugation together preserve the zero set under these maps, but invariance of a set under a reflection does not imply that each of its points is fixed.

For example, for $a,b>0$,
\begin{equation}
 P(z)=\bigl((z-a)^2+b^2\bigr)\bigl((z+a)^2+b^2\bigr)
 \label{eq:quartet}
\end{equation}
is even and real on the real axis, while its zeros are $\pm a\pm ib$. Taking $0<b<1/2$ places the corresponding $s=1/2+iz$ zeros inside the critical strip but off its central line. This elementary example is not a counterexample to RH: it has no zeta Euler product. It refutes only the inference from reflection symmetry to pointwise zero confinement.

\subsection{The half-density calculation}
There is nevertheless a precise analytical reason for a factor $1/2$ in a scale representation. On $L^2(\R_+,\dd x)$ define, for $a>0$,
\begin{equation}
 (D_af)(x)=a^{1/2}f(ax).
 \label{eq:dilation}
\end{equation}
A change of variables gives $\|D_af\|_2=\|f\|_2$. Under the logarithmic unitary map
\begin{equation}
 (Uf)(r)=\e^{r/2}f(\e^r),\qquad
 U:L^2(\R_+,\dd x)\longrightarrow L^2(\R,\dd r),
 \label{eq:logunitary}
\end{equation}
dilation becomes translation: $UD_{\e^v}U^{-1}g(r)=g(r+v)$.

\begin{proposition}[Half-density generator]
\label{prop:half}
The operator
\begin{equation}
 B=-i\left(x\frac{\dd}{\dd x}+\frac12\right)
 \label{eq:dilationgenerator}
\end{equation}
on $C_c^\infty(\R_+)$ is essentially self-adjoint. Its closure is unitarily equivalent to $-i\dd/\dd r$ on $L^2(\R)$ and has purely continuous spectrum $\R$.
\end{proposition}
\begin{proof}
Differentiating Equation~\eqref{eq:logunitary} gives $UBU^{-1}=-i\dd/\dd r$ on $C_c^\infty(\R)$. The Fourier transform realizes the closure of the latter as multiplication by a real variable, with domain consisting of functions whose Fourier transforms remain square-integrable after multiplication by that variable. This multiplication operator is self-adjoint and has no square-integrable point eigenvectors. Its spectrum is $\R$. Unitary equivalence gives the assertion.
\end{proof}

The Mellin transform of $f$ on $s=1/2+it$ becomes a Fourier transform of $Uf$:
\begin{equation}
 \int_0^\infty f(x)x^{-1/2+it}\dd x
 =\int_{-\infty}^{\infty}(Uf)(r)\e^{itr}\dd r.
 \label{eq:mellinhalf}
\end{equation}
The real part $1/2$ is fixed here by the Lebesgue measure $\dd x$ and its half-density correction. It is not obtained by subtracting three source dimensions from four or by treating a temporal coordinate as imaginary. The choice $L^2(\R_+,\dd x/x)$ shifts the corresponding normalization, illustrating why the measure must be declared before assigning geometric significance to the exponent.

This calculation gives a real spectral axis but not the arithmetic spectrum. The dilation operator has continuous spectrum, whereas $F$ has an infinite discrete zero set. Discreteness, multiplicity, and prime-sensitive trace data remain to be supplied. The calculation is thus a useful local geometric ingredient and not a proof of RH. Its role parallels the distinction between a formal dilation Hamiltonian and a fully specified spectral realization in the Berry--Keating programme \cite{berrykeating}.

\section{Positive theta geometry and the nonlinear moment step}
\label{sec:theta}
\subsection{An exact positive-kernel representation}
Let
\begin{equation}
 \vartheta_+(x)=\sum_{n=1}^{\infty}\e^{-\pi n^2x},\qquad
 h(u)=\e^{u/2}\vartheta_+(\e^{2u}),\quad u\geq0.
\end{equation}
The theta transformation underlying the functional equation gives the entire identity
\begin{equation}
 \xi(s)=\frac12+\frac{s(s-1)}2
 \int_1^\infty\vartheta_+(x)
 \left(x^{s/2}+x^{(1-s)/2}\right)\frac{\dd x}{x}.
 \label{eq:thetaxi}
\end{equation}
This is a standard form of Riemann's theta representation \cite{riemann,titchmarsh,dlmf}. With $s=1/2+iz$ and $x=\e^{2u}$, it becomes
\begin{equation}
 \Xi(z)=\frac12-2(z^2+\tfrac14)
                    \int_0^\infty h(u)\cos(zu)\dd u.
 \label{eq:thetaint}
\end{equation}

The full theta relation $\vartheta(x)=x^{-1/2}\vartheta(1/x)$, where $\vartheta=1+2\vartheta_+$, implies $\vartheta'(1)=-\vartheta(1)/4$. Therefore
\begin{equation}
 h'(0)=\frac12\vartheta_+(1)+2\vartheta_+'(1)=-\frac14.
\end{equation}
Twice integrating by parts in Equation~\eqref{eq:thetaint}, with all terms at infinity vanishing, yields
\begin{equation}
 \Xi(z)=\int_0^\infty\Phi(u)\cos(zu)\dd u,
 \qquad \Phi(u)=2\bigl(h''(u)-\tfrac14h(u)\bigr).
 \label{eq:positivekernel}
\end{equation}
Termwise differentiation gives the explicit formula
\begin{equation}
 \Phi(u)=4\sum_{n=1}^{\infty}
 \left(2\pi^2n^4\e^{9u/2}-3\pi n^2\e^{5u/2}\right)
 \e^{-\pi n^2\e^{2u}}.
 \label{eq:phiseries}
\end{equation}
Every summand is strictly positive for $u\geq0$, because $2\pi n^2\e^{2u}>3$. The double-exponential decay justifies differentiation and integration on compact subsets of the complex $z$-plane. In particular, Equation~\eqref{eq:positivekernel} establishes the real-axis positivity used in Lemma~\ref{lem:product}.

\subsection{Raw moments and logarithmic moments are different objects}
Define
\begin{equation}
 A_{2j}=\int_0^\infty\Phi(u)u^{2j}\dd u,
 \qquad F(w)=\sum_{j=0}^{\infty}f_jw^j,
 \qquad f_j=\frac{(-1)^jA_{2j}}{(2j)!A_0}.
 \label{eq:rawmoments}
\end{equation}
The $A_{2j}$ are moments of an explicitly positive measure. Their ordinary moment matrices are positive. The arithmetic matrices needed for RH, however, are generated by a logarithmic derivative:
\begin{equation}
 M(w)=-\frac{F'(w)}{F(w)}=\sum_{n=0}^{\infty}m_nw^n
 \quad\text{near }w=0.
 \label{eq:logmoments}
\end{equation}
Equating coefficients in $MF=-F'$ gives
\begin{equation}
 m_n=-(n+1)f_{n+1}-\sum_{k=1}^{n}f_km_{n-k}.
 \label{eq:recurrence}
\end{equation}
For example,
\begin{align}
 m_0&=-f_1, &m_1&=f_1^2-2f_2,\\
 m_2&=-f_1^3+3f_1f_2-3f_3.
 \label{eq:firstmoments}
\end{align}
This nonlinear transformation is exactly where positivity ceases to be automatic. The positivity of $\Phi$ supplies a Hilbert norm for raw polynomial moments in $u$, but it does not yet supply a Hilbert norm for the $m_n$. An ECL interpretation must address the latter norm, not substitute the former.

\subsection{Why positivity of a Fourier kernel is insufficient}
Let $g$ be the centered Gaussian probability density whose Fourier transform is $\e^{-z^2}$. The strictly positive even function
\begin{equation}
 k(u)=2g(u)+\frac12g(u-1)+\frac12g(u+1)
\end{equation}
has entire Fourier transform
\begin{equation}
 \widehat k(z)=\e^{-z^2}(2+\cos z).
\end{equation}
Its zeros include $z=\pi\pm i\arcosh 2$, which are nonreal. Thus an even positive kernel and an entire transform do not suffice for real-rootedness. This example does not have the exact order or arithmetic structure of $\Xi$, and it is not offered as a counterexample inside the zeta class. It isolates the invalid positivity implication.

The geometric task is consequently sharper than producing a positive density. One must show that the \emph{logarithmic} moment form induced by the actual theta kernel is positive in every polynomial degree. The following section makes that statement equivalent to the original hypothesis, without assuming the zeros are simple or computing their ordinates.

\section{The positive moment criterion}
\label{sec:moments}
\subsection{An arithmetic quadratic form}
For $N\geq0$, define the real symmetric Hankel matrix
\begin{equation}
 H_N=(m_{i+j})_{i,j=0}^{N}.
 \label{eq:hankel}
\end{equation}
For a polynomial $p(t)=\sum_{j=0}^{N}c_jt^j$, set
\begin{equation}
 q_\xi(p)=\sum_{i,j=0}^{N}\overline{c_i}c_jm_{i+j}.
 \label{eq:qxi}
\end{equation}
The subscript records that the form is extracted from the specific Riemann completion, not selected from an arbitrary positive model. We use inner products conjugate-linear in the first variable throughout. For $c=a+ib$ with real vectors $a,b$, symmetry gives $c^*H_Nc=a^TH_Na+b^TH_Nb$; thus positivity over real coefficients is equivalent to positivity over complex coefficients.

The canonical product gives an exact arithmetic interpretation of its coefficients:
\begin{equation}
 M(w)=\sum_j\frac{1}{w_j-w},\qquad
 m_n=\sum_jw_j^{-n-1}.
 \label{eq:powersums}
\end{equation}
These series converge absolutely near the origin, and the first series converges normally on compact sets avoiding the zeros. Conjugate pairing makes every $m_n$ real even without RH. With a sum over \emph{all} nontrivial zeros, rather than the upper-half-plane representatives used here, the corresponding reciprocal-square power sums are $2m_n$ \cite{voros}.

\begin{lemma}[Analytic moment growth forces compact support]
\label{lem:compactmoment}
Suppose $(m_n)_{n\geq0}$ is a real Hamburger moment sequence and $\sum_{n\geq0}m_nw^n$ has positive radius of convergence. Then every positive representing measure has compact support.
\end{lemma}
\begin{proof}
Choose $r>0$ inside the radius of convergence. Cauchy's estimate supplies a constant $C_r$ with $|m_n|\leq C_rr^{-n}$. Let $\mu$ be a positive representing measure. If $\mu(\{|t|\geq A\})=c>0$ for some $A>r^{-1}$, then
\[
 m_{2n}=\int_\R t^{2n}\dd\mu(t)\geq cA^{2n}.
\]
This contradicts $m_{2n}\leq C_rr^{-2n}$ as $n\to\infty$. Hence $\supp\mu\subseteq[-r^{-1},r^{-1}]$.
\end{proof}

\begin{theorem}[Reciprocal-square Hamburger criterion]
\label{thm:moment}
For the function $F$ in Equation~\eqref{eq:normal}, the following are equivalent:
\begin{enumerate}
 \item the Riemann hypothesis;
 \item $H_N$ is positive semidefinite for every $N\geq0$;
 \item $H_N$ is positive definite for every $N\geq0$;
 \item $\det H_N>0$ for every $N\geq0$.
\end{enumerate}
No simplicity assumption on the zeros is required.
\end{theorem}
\begin{proof}
Assume RH. Write the positive ordinates, with multiplicity, as $\gamma_j$ and put $\lambda_j=\gamma_j^{-2}$. Lemma~\ref{lem:product} gives $\sum_j\lambda_j<\infty$. Define the finite positive measure
\begin{equation}
 \mu_\xi=\sum_j\lambda_j\delta_{\lambda_j}.
 \label{eq:selfweighted}
\end{equation}
Then $m_n=\int t^n\dd\mu_\xi(t)$, so
\begin{equation}
 q_\xi(p)=\int|p(t)|^2\dd\mu_\xi(t)
         =\sum_j\lambda_j|p(\lambda_j)|^2\geq0.
 \label{eq:momentnorm}
\end{equation}
There are infinitely many distinct $\lambda_j$, since $\xi$ has infinitely many zeros and each zero has finite multiplicity. A nonzero polynomial cannot vanish on all of them. Thus Equation~\eqref{eq:momentnorm} is strictly positive for nonzero $p$, proving statement 3.

Conversely, assume statement 2. The Hamburger moment theorem gives a positive Borel measure $\mu$ on $\R$ satisfying $m_n=\int t^n\dd\mu(t)$ \cite{schmudgen}. Lemma~\ref{lem:compactmoment} makes its support compact. Therefore
\begin{equation}
 G_\mu(w)=\int_\R\frac{\dd\mu(t)}{1-wt}
 \label{eq:momenttransform}
\end{equation}
is holomorphic in each open nonreal half-plane and agrees with $M(w)$ near $w=0$. Indeed, on a sufficiently small disc the geometric series can be integrated uniformly. The identity theorem for meromorphic functions then gives $G_\mu=-F'/F$ throughout each nonreal half-plane.

If $F$ had a nonreal zero $w_*$ of multiplicity $d\geq1$, its logarithmic derivative would have the nonremovable local term
\[
 -\frac{F'(w)}{F(w)}=-\frac{d}{w-w_*}+\text{holomorphic part}.
\]
This contradicts the holomorphy of Equation~\eqref{eq:momenttransform}. Consequently all zeros of $F$ are real. Lemma~\ref{lem:product} excludes $(-\infty,0]$, so all its zeros are positive. Corollary~\ref{cor:positivezeros} proves RH.

Statement 3 implies statement 2. Statement 3 also implies statement 4. Conversely, statement 4 gives positivity of every leading principal minor of each finite $H_N$, and Sylvester's criterion gives statement 3. This completes the equivalence.
\end{proof}

Theorem~\ref{thm:moment} is a reciprocal-square formulation within the classical real-zero/moment method \cite{grommer,schmudgen,bariczstampach}. The proof above supplies the analytic normalization and compact-support argument needed here; it does not claim priority for a new real-zero criterion. The relation to Grommer's inverse-power sequence is derived in Section~\ref{subsec:grommernormalization}. Its utility for TCGS--SEQUENTION is that the target is now a definite positive form, Equation~\eqref{eq:qxi}, to which a readout norm can be compared.

\begin{remark}[A distinction about minors]
It is important that statement 4 uses \emph{strictly positive} leading determinants. Merely requiring all leading determinants to be nonnegative is not, in general, a sufficient test for positive semidefiniteness of a real symmetric matrix. The unambiguous non-strict formulation is statement 2, which tests each complete matrix. No argument below substitutes one condition for the other.
\end{remark}

\subsection{What the first conditions actually test}
The first two matrix conditions are
\begin{equation}
 m_0\geq0,\qquad m_0m_2-m_1^2\geq0.
 \label{eq:lowconditions}
\end{equation}
Positivity of each individual $m_n$ is weaker than the Hankel requirement. The second inequality already compares three logarithmic moments. In a positive realization it is a Cauchy--Schwarz inequality between the constant polynomial and the coordinate polynomial. In arithmetic form it is a nonlinear relation among the first three nonconstant coefficients of $F$.

This illustrates the exact conceptual gain. The question is not whether a list of scalar values looks positive, but whether all finite combinations define a compatible Hilbert geometry. Every degree enlarges the class of combinations that must be controlled. A positive measure for the raw theta moments does not automatically establish these logarithmic-moment inequalities, because Equation~\eqref{eq:recurrence} does not preserve moment positivity by definition.


\subsection{The classical Grommer criterion and the reciprocal-square indexing}
\label{subsec:grommernormalization}
Grommer's real-zero criterion is stated in terms of Hankel determinants of inverse powers of zeros \cite{grommer}. We use the order-one formulation given explicitly in \cite[Theorem~10]{bariczstampach}: for the class relevant here, the moments begin at the inverse \emph{second} power, and strict positivity of the complete determinant sequence characterizes reality of all zeros. The Riemann function $\Xi$ has real Taylor coefficients, order one, no zero at the origin, infinitely many distinct zeros, and the needed inverse-square summability. These hypotheses are important; strict positivity is not asserted for a finite-support zero set. General background on real entire functions is provided by Levin \cite{levin}.

Let $(z_\ell)$ enumerate \emph{all} zeros of $\Xi$, with multiplicity, and define
\begin{equation}
 g_k=\sum_\ell z_\ell^{-k-2},\qquad
 \mathsf G_n=(g_{i+j})_{i,j=0}^{n-1},\qquad n\geq1.
 \label{eq:grommermoments}
\end{equation}
Absolute convergence and the pairing $z\leftrightarrow-z$ give the exact identities
\begin{equation}
 g_{2j}=2m_j,\qquad g_{2j+1}=0.
 \label{eq:grommerparity}
\end{equation}
Introduce the shifted Hankel matrix
\begin{equation}
 H_N^{(+)}=(m_{i+j+1})_{i,j=0}^{N}.
 \label{eq:shiftedhankel}
\end{equation}
Permuting the rows and columns of $\mathsf G_n$ into even and odd index order therefore yields
\begin{equation}
 \mathsf P_n\mathsf G_n\mathsf P_n^{*}
 =2H_{\lceil n/2\rceil-1}\oplus2H_{\lfloor n/2\rfloor-1}^{(+)},
 \label{eq:grommerblocks}
\end{equation}
where an index $-1$ denotes an absent block. Equivalently, with the determinant of the empty matrix equal to one,
\begin{equation}
 \det\mathsf G_n
 =2^n\det H_{\lceil n/2\rceil-1}
       \det H_{\lfloor n/2\rfloor-1}^{(+)}.
 \label{eq:grommerdeterminants}
\end{equation}
These formulas identify the normalization and the missing shifted block explicitly. The single tower in Theorem~\ref{thm:moment} is not literally the full Grommer matrix sequence. Its sufficiency uses additional structure of the fixed function $F$: genus-zero factorization, analyticity of the logarithmic moment transform, and the absence of zeros on $(-\infty,0]$. The compact-support/pole argument proves positivity of its zeros from the unshifted tower; then Equation~\eqref{eq:selfweighted} gives both towers as positive moment matrices. For an arbitrary Hamburger sequence, unshifted positivity alone does not imply positivity of the shifted tower. For a Hamburger sequence, the additional condition $H_N^{(+)}\succeq0$ for all $N$ is precisely the Stieltjes support condition for a representing measure on $[0,\infty)$ \cite{schmudgen}. Thus, for this particular $F$, Theorem~\ref{thm:moment} upgrades Hamburger positivity to Stieltjes positivity because Lemma~\ref{lem:product} has already excluded $(-\infty,0]$ from the zero set. No simplicity assumption is introduced by this comparison. The related Laguerre and Tur\'an inequalities \cite{cravencsordas} are not substituted for this all-degree Hankel requirement.

\section{Fredholm determinants and a prime-side closure identity}
\label{sec:trace}
\subsection{Positive trace-class realization}
For a positive trace-class operator $K$ on a separable Hilbert space, the Fredholm determinant is
\begin{equation}
 D_K(w)=\det(\Id-wK)=\prod_j(1-w\lambda_j),
 \label{eq:fredholm}
\end{equation}
where $(\lambda_j)$ are the positive eigenvalues with multiplicity. A kernel of $K$ contributes no factor. The product converges locally uniformly because $\sum_j\lambda_j<\infty$, and its zeros are positive real \cite{simon}.

\begin{theorem}[Determinant and trace equivalences]
\label{thm:determinant}
The following statements are equivalent:
\begin{enumerate}
 \item RH holds;
 \item there exists a positive trace-class operator $K$ such that $F(w)=D_K(w)$ for every $w\in\C$;
 \item there exists a positive trace-class operator $K$ such that
 \begin{equation}
   \Tr K^{n+1}=m_n\qquad(n\geq0).
   \label{eq:tracemoments}
 \end{equation}
\end{enumerate}
\end{theorem}
\begin{proof}
Under RH, take $K$ on $\ell^2$ to be diagonal with eigenvalues $\gamma_j^{-2}$. Lemma~\ref{lem:product} proves trace-class membership and the determinant identity. This gives statement 2, and differentiating the locally convergent product gives statement 3.

Conversely, if statement 3 holds, then for $|w|<\|K\|^{-1}$,
\begin{equation}
 -\frac{D_K'(w)}{D_K(w)}
 =\Tr\bigl(K(\Id-wK)^{-1}\bigr)
 =\sum_{n=0}^{\infty}\Tr K^{n+1}w^n=M(w).
 \label{eq:logdet}
\end{equation}
Both determinants are normalized to one at the origin. Equality of logarithmic derivatives therefore gives $D_K=F$ near the origin and, by analyticity, on $\C$. Statement 2 implies RH because every zero of a positive Fredholm determinant lies on $(0,\infty)$.
\end{proof}

The diagonal operator in the forward implication is a representation \emph{conditional on RH}. It is not a source-derived construction and cannot be inserted as the first step in a proof. A genuine construction must specify $K$, including its measure, domain, boundary conditions, and normalization, from data independent of the assumed location of the arithmetic zeros. Theorem~\ref{thm:determinant} then supplies an exact test of that construction.

\subsection{Where the primes enter}
Let $\Lambda(n)$ be the von Mangoldt function: $\Lambda(p^k)=\log p$ for primes $p$ and integers $k\geq1$, and $\Lambda(n)=0$ otherwise. In the half-plane $\Re s>1$, differentiating the absolutely convergent Euler product gives
\begin{equation}
 -\frac{\zeta'(s)}{\zeta(s)}=\sum_{n=2}^{\infty}\frac{\Lambda(n)}{n^s}.
 \label{eq:mangoldt}
\end{equation}
Consequently, with $\psi_\Gamma=\Gamma'/\Gamma$,
\begin{equation}
 \frac{\xi'(s)}{\xi(s)}
 =\frac1s+\frac1{s-1}-\frac12\log\pi
  +\frac12\psi_\Gamma(s/2)
  -\sum_{n=2}^{\infty}\frac{\Lambda(n)}{n^s}.
 \label{eq:arithlog}
\end{equation}
Both identities are standard arithmetic consequences of Equation~\eqref{eq:euler} \cite{titchmarsh}. The subscript on $\psi_\Gamma$ prevents confusion with the Chebyshev prime-counting function.

\begin{theorem}[Exact prime-trace closure]
\label{thm:primetrace}
RH is equivalent to the existence of a positive trace-class operator $K$ satisfying, for every $s$ in some nonempty open real interval $J\subset(1,\infty)$,
\begin{align}
 2(s-\tfrac12)\Tr\!\left[K\bigl(\Id+(s-\tfrac12)^2K\bigr)^{-1}\right]
 &=\frac1s+\frac1{s-1}-\frac12\log\pi
   +\frac12\psi_\Gamma(s/2)\notag\\[-0.2em]
 &\hspace{0.7em}-\sum_{n=2}^{\infty}\frac{\Lambda(n)}{n^s}.
 \label{eq:primebridge}
\end{align}
If the identity holds on such an interval, it holds for every real $s>1$.
\end{theorem}
\begin{proof}
Assume the stated operator exists. Put $t=s-1/2$ and define the entire function
\[
 R_K(s)=D_K\bigl(-(s-\tfrac12)^2\bigr).
\]
For real $s>1$, $R_K(s)>0$, and Equation~\eqref{eq:logdet} gives
\[
 \frac{R_K'(s)}{R_K(s)}
 =2t\Tr\bigl[K(\Id+t^2K)^{-1}\bigr].
\]
Equations~\eqref{eq:arithlog} and \eqref{eq:primebridge} imply $R_K(s)=c\xi(s)$ on $J$, where $c$ is constant. The identity theorem extends this equality to all complex $s$. Evaluating at $s=1/2$ gives $c=\xi(1/2)^{-1}$. Therefore $D_K=F$ on the negative real axis and hence everywhere. Theorem~\ref{thm:determinant} proves RH. Conversely, the diagonal realization under RH satisfies Equation~\eqref{eq:logdet}, and differentiating $D_K(-(s-1/2)^2)=\xi(s)/\xi(1/2)$ yields Equation~\eqref{eq:primebridge} for all real $s>1$.
\end{proof}

Equation~\eqref{eq:primebridge} is the decisive arithmetic identification in this paper. Its left side is positive-operator geometry. Its right side is fixed by prime powers, the gamma factor, and the elementary completion factors. Neither side contains adjustable coefficients chosen from zero ordinates. The identity makes precise what must be achieved beyond symmetry, leading asymptotics, or a generic self-adjoint construction.

The theorem also separates analytical continuation from divergent manipulation. The arithmetic side is evaluated only where the prime-power series converges absolutely. Equality of entire functions is obtained afterward by the identity theorem. Nothing requires multiplying the Euler factors on $\Re s=1/2$, interchanging a divergent prime sum with a trace, or assuming the desired zero-free region in order to define the operator.

\subsection{Why leading density cannot replace the identity}
An eigenvalue counting law determines only an averaged spectral distribution. The full resolvent trace in Equation~\eqref{eq:primebridge} determines much more: all inverse-power traces near the origin of the determinant, every zero with its multiplicity, and the normalization of the entire function. Two operators can have the same first several counting terms and still have different determinants. Section~\ref{sec:exponential} supplies an explicit example at the precise zeta normalization.

The arithmetic interpretation is correspondingly exact. Under RH, the classical error estimate is $\pi(x)-\Li(x)=O(\sqrt{x}\log x)$, with $\Li$ interpreted in its usual principal-value normalization \cite{bombieri}. The relevant fine-scale comparison is with $\Li(x)$, not simply with $x/\log x$. The trace identity retains the complete prime-power information required by this connection, whereas matching the smooth counting term alone cannot recover it.

\section{Transferring constitutive positivity to arithmetic}
\label{sec:gram}
\subsection{The exact bridge}
The positive operator $O_\sigma=L_\sigma^*L_\sigma$ from Equation~\eqref{eq:observability} suggests a precise sufficient condition. The condition is not an additional assertion that all arithmetic is physically generated by a Hilbert space. It specifies what an arithmetic representation of a particular source reduction would have to preserve.

\begin{assumption}[Arithmetic Gram identification]
\label{ass:gram}
For an admissible selector $\sigma$, there exists a complex-linear map
\begin{equation}
 J_\sigma:\C[t]\longrightarrow\Dom(L_\sigma)
\end{equation}
such that, for all polynomials $p,q$,
\begin{equation}
 \left\langle L_\sigma J_\sigma p,L_\sigma J_\sigma q\right\rangle
 =\sum_{i,j}\overline{p_i}q_jm_{i+j}.
 \label{eq:grambridge}
\end{equation}
Here $p_i,q_j\in\C$ are polynomial coefficients, and the bar on $p_i$ denotes complex conjugation. Both sides are conjugate-linear in $p$ and linear in $q$; in particular, replacing $p$ by $ip$ multiplies each side by $-i$. The source/readout norms and the map $J_\sigma$ are fixed independently of any assumption that the zeta zeros lie on the critical line.
\end{assumption}

\begin{theorem}[Constitutive Gram transfer]
\label{thm:gramtransfer}
A realization satisfying Assumption~\ref{ass:gram} proves RH. Alternatively, it is sufficient to construct a positive trace-class reduction $K$ of the declared source/readout representation and prove Equation~\eqref{eq:tracemoments} in every degree.
\end{theorem}
\begin{proof}
Setting $p=q$ in Equation~\eqref{eq:grambridge} gives $q_\xi(p)=\|L_\sigma J_\sigma p\|^2\geq0$. Hence every $H_N$ is positive semidefinite and Theorem~\ref{thm:moment} applies. The trace formulation follows from Theorem~\ref{thm:determinant}.
\end{proof}

The proof is short because the difficult identification is displayed rather than hidden. The framework supplies a legitimate class of positive readout norms; it does not, merely by doing so, prove Equation~\eqref{eq:grambridge}. Constructing $J_\sigma$ by first assuming that $q_\xi$ is positive would reverse the intended implication. Likewise, choosing a Hilbert norm so that an already selected arithmetic matrix becomes a Gram matrix would encode the unproved target in the input \cite{tcgsoperator}.

For a positive trace-class $K$, the vectors $K^{n+1/2}$ are Hilbert--Schmidt operators and satisfy
\begin{equation}
 \left\langle K^{i+1/2},K^{j+1/2}\right\rangle_{\mathrm{HS}}
 =\Tr K^{i+j+1}.
 \label{eq:hsgram}
\end{equation}
Thus a successful trace construction produces the required polynomial Gram geometry explicitly. Equation~\eqref{eq:hsgram} is useful in the forward direction, after $K$ and its arithmetic trace have been derived. It is not a way to bypass the derivation of $K$.

\subsection{Selector covariance}
The framework organizes physically equivalent selectors by comparison arrows, not by the assertion that every boundary condition describes the same physical readout \cite{tcgsfoundation,tcgsoperator}. Suppose an arrow $\kappa:\sigma\to\sigma'$ has unitary input and output maps $U_\kappa,V_\kappa$ with the required domain transport, so that
\begin{equation}
 L_{\sigma'}=V_\kappa L_\sigma U_\kappa^{-1},\qquad
 J_{\sigma'}=U_\kappa J_\sigma.
 \label{eq:naturality}
\end{equation}

\begin{proposition}[Naturality of arithmetic positivity]
\label{prop:natural}
Under Equation~\eqref{eq:naturality}, an exact Gram identification for $\sigma$ is an exact Gram identification for $\sigma'$.
\end{proposition}
\begin{proof}
For any $p,q$, substitute Equation~\eqref{eq:naturality} into the left side of Equation~\eqref{eq:grambridge}. The input unitaries cancel, and the common output unitary preserves the inner product. The resulting expression equals its value for $\sigma$.
\end{proof}

This is the appropriate invariant statement: the arithmetic form is preserved along the declared equivalence arrows. Changing to an inequivalent boundary phase or selecting a different potential is not automatically such an arrow. Requiring those distinct models to share one source rule is a transfer requirement; it is not a theorem that their determinants coincide. This distinction becomes concrete when the Dirichlet exponential model is tested below.

\subsection{An infinite-dimensional requirement}
\begin{theorem}[Finite-rank obstruction]
\label{thm:rank}
No map satisfying Equation~\eqref{eq:grambridge} can have finite-dimensional image under $L_\sigma J_\sigma$. In particular, one fixed finite-dimensional protected response sector cannot realize the full arithmetic moment tower.
\end{theorem}
\begin{proof}
Suppose $\rank(L_\sigma J_\sigma)=r<\infty$. For $N\geq r$, the $N+1$ vectors $L_\sigma J_\sigma t^j$, $0\leq j\leq N$, have a singular Gram matrix. Thus $\rank H_N\leq r$. But Equation~\eqref{eq:grambridge} implies RH by Theorem~\ref{thm:gramtransfer}, and RH implies $H_N$ is positive definite by Theorem~\ref{thm:moment}. This contradicts $\rank H_N<N+1$.
\end{proof}

The obstruction concerns representation rank, not source dimension. Function spaces over a four-dimensional manifold are normally infinite-dimensional. The theorem therefore excludes a finite arithmetic truncation masquerading as a complete representation, not a four-dimensional source. Finite sectors remain admissible for finite diagnostics, provided their rank limit and truncation error are reported rather than promoted into an all-degree proof.

\subsection{The missing construction in operational form}
An executable source route must determine three objects in the correct order: the source/readout Hilbert realization, an arithmetic insertion map such as $J_\sigma$, and an identity proving that its Gram coefficients equal the $m_n$ extracted from $\xi$. A determinant route additionally needs a trace-class positive reduction $K$, with a declared relation to $O_\sigma$. There is no automatic rule $K=O_\sigma$, since an observability operator can be unbounded, have continuous spectrum, or fail every required trace condition \cite{tcgsoperator}.

The strongest compact statement is therefore
\begin{equation}
 \boxed{\text{source-derived positive representation}
 \; +\;\text{exact arithmetic identification}
 \;\Longrightarrow\;\mathrm{RH}.}
 \label{eq:typedimplication}
\end{equation}
The two premises are mathematically different. The remainder of the paper tests the extent to which a concrete geometric model can satisfy both.

\subsection{A positive compatible-cochain sector}
\label{subsec:cochainpositive}
A concrete algebraic starting point is available without prescribing a one-dimensional potential. On the compatible cochain complex of Equation~\eqref{eq:compatiblecochains}, fix an admissible reference phase and a positive weighted Hodge matrix $M_{2,\chi,r}$. With any Gate-T period sector held fixed, variation of the exact reciprocal component gives
\begin{equation}
 C^{\mathrm{rec}}_r=d_{1,r}^{*}M_{2,\chi,r}d_{1,r},
 \qquad a^{*}C^{\mathrm{rec}}_ra
 =(d_{1,r}a)^{*}M_{2,\chi,r}(d_{1,r}a)\geq0.
 \label{eq:cochainstiffness}
\end{equation}
Here the star on a coefficient matrix denotes conjugate transpose; it is not an undeclared metric adjoint. If $M_{1,r}>0$ is the one-cochain mass matrix, the Euclidean-normalized realization is $M_{1,r}^{-1/2}C^{\mathrm{rec}}_rM_{1,r}^{-1/2}$. Its null space, period constraints, and boundary conditions must be treated explicitly. Only independently established redundancies may be removed \cite{tcgsconservative,tcgsoperator}.

Equation~\eqref{eq:cochainstiffness} is a positive sector, not a positivity theorem for the full coupled foam Hessian. The phase potential, field couplings, singular boundary term, and admissible variation space must all be included before a total second variation is assessed. Nonnegative energy values do not imply a nonnegative Hessian at every reference configuration. Moreover, an energy Hessian is distinct from the physical-readout derivative construction in Equation~\eqref{eq:observability} \cite{tcgsfoam,tcgsoperator}.

\subsection{Positive boundary response by interior elimination}
\label{subsec:schur}
Suppose the declared quadratic realization at refinement $r$, after its stated constraints, is a Hermitian matrix $\mathsf C_r\geq0$. Partition its variables into interior and retained boundary/interface coefficients:
\begin{equation}
 \mathsf C_r=
 \begin{pmatrix}C_{ii,r}&C_{ib,r}\\C_{bi,r}&C_{bb,r}\end{pmatrix},
 \qquad C_{ii,r}>0,\quad C_{bi,r}=C_{ib,r}^{*}.
 \label{eq:sourceblock}
\end{equation}
The boundary contribution in Equation~\eqref{eq:sharedboundaryvariation}, when constructed, belongs in this same quadratic assembly. It is not added independently after arithmetic comparison. The following result is standard Schur-complement algebra; its role is to identify a positivity-preserving reduction compatible with the declared source construction, not to claim a new matrix theorem. Graph-theoretic Kron reduction supplies a familiar realization of this mechanism \cite{kron}.

\begin{proposition}[Positive interface reduction]
\label{prop:schurpositive}
For Equation~\eqref{eq:sourceblock}, define
\begin{equation}
 \mathsf S_r=C_{bb,r}-C_{bi,r}C_{ii,r}^{-1}C_{ib,r},
 \qquad
 \mathsf T_r b=\begin{pmatrix}-C_{ii,r}^{-1}C_{ib,r}b\\b\end{pmatrix}.
 \label{eq:schurresponse}
\end{equation}
Then $\mathsf S_r=\mathsf T_r^{*}\mathsf C_r\mathsf T_r\geq0$, and
\begin{equation}
 b^{*}\mathsf S_rb
 =\min_{a}\begin{pmatrix}a\\b\end{pmatrix}^{*}
                  \mathsf C_r\begin{pmatrix}a\\b\end{pmatrix}.
 \label{eq:schurminimum}
\end{equation}
If $\mathsf C_r>0$, then $\mathsf S_r>0$.
\end{proposition}
\begin{proof}
Complete the square:
\begin{align*}
 \begin{pmatrix}a\\b\end{pmatrix}^{*}\mathsf C_r
 \begin{pmatrix}a\\b\end{pmatrix}
 &=\left(a+C_{ii,r}^{-1}C_{ib,r}b\right)^{*}C_{ii,r}
   \left(a+C_{ii,r}^{-1}C_{ib,r}b\right)+b^{*}\mathsf S_rb.
\end{align*}
The first term has its unique minimum zero at the stated interior extension. Substitution gives both the minimum identity and the congruence. If the full matrix is positive definite, the extension of any nonzero $b$ is nonzero and has strictly positive quadratic value.
\end{proof}

With a declared positive boundary mass matrix $M_{b,r}$, put
\begin{equation}
 B_r=M_{b,r}^{-1/2}\mathsf S_rM_{b,r}^{-1/2},
 \qquad K_r=B_r^{-1}
 \label{eq:boundarycandidate}
\end{equation}
when $B_r>0$ on the declared retained space. At each finite refinement this is a positive trace-class candidate. A continuum trace-class limit, dimensional normalization, and arithmetic identity remain independent requirements. A zero eigenvalue may not be removed, shifted, or assigned a pseudoinverse solely to force agreement with $F$; any quotient, restriction, or shift changes the candidate and must be fixed by its source and readout definition.

The inverse interior block in Equation~\eqref{eq:schurresponse} can generate a dense boundary response from a sparse local assembly. This supplies a calculable route to nonlocality rather than an independently fitted nonlocal kernel. The boundary is an auxiliary interface of the mathematical representation. It is not automatically the physical Shadow or a complete holographic encoding. The separate selector and physical-registration maps remain required \cite{tcgsconservative,tcgsfoundation}.

\begin{proposition}[Selector covariance of the reduced candidate]
\label{prop:schurcovariance}
Suppose equivalent selectors transport Equation~\eqref{eq:sourceblock} by $U_i\oplus U_b$, with $U_i,U_b$ unitary, and transport the boundary mass by $M_b'=U_bM_bU_b^{*}$. Then
\begin{equation}
 \mathsf S'=U_b\mathsf S U_b^{*},\qquad
 B'=U_bBU_b^{*},\qquad K'=U_bKU_b^{*}
 \label{eq:schurnaturality}
\end{equation}
whenever the inverses exist. Their determinants and power traces agree.
\end{proposition}
\begin{proof}
Insert the blockwise conjugation into Equation~\eqref{eq:schurresponse}; the interior unitaries cancel through $C_{ii}^{-1}$. Functional calculus transports the positive square root of $M_b$, and inversion gives the last relation. Trace and determinant are unitary invariants.
\end{proof}

The block-compatible transport is a hypothesis to be checked, not a property of arbitrary selector changes. In the continuum, one may instead define a boundary form by minimizing a source form over extensions with prescribed trace. A densely defined closed nonnegative boundary form has an associated positive self-adjoint operator; a strictly positive lower bound and a trace-class inverse require further estimates \cite{reed}. Neither closability nor spectral compactness is supplied by formal elimination alone.

\subsection{Static response is not the complete bulk determinant}
\label{subsec:spectralelimination}
The arithmetic object must be fixed before a spectral comparison. A static response operator and a parameter-dependent elimination of the complete source problem are different objects. Let $\mathsf C_r>0$ and $\mathsf M_r>0$, and define the full spectral pencil $\mathsf P_r(w)=\mathsf C_r-w\mathsf M_r$. Wherever its interior block is invertible, write
\begin{equation}
 \mathsf S_r(w)=P_{bb,r}(w)
 -P_{bi,r}(w)P_{ii,r}(w)^{-1}P_{ib,r}(w).
 \label{eq:spectralschur}
\end{equation}
Block elimination gives the exact factorization
\begin{equation}
 \frac{\det(\mathsf C_r-w\mathsf M_r)}{\det\mathsf C_r}
 =\frac{\det P_{ii,r}(w)}{\det C_{ii,r}}
  \frac{\det\mathsf S_r(w)}{\det\mathsf S_r(0)}.
 \label{eq:bulkboundarydeterminant}
\end{equation}
To verify it, left-multiply $\mathsf P_r(w)$ by the unit-determinant block triangular matrix that subtracts $P_{bi,r}(w)P_{ii,r}(w)^{-1}$ times its first block row from the second. The result is block triangular. Normalization uses the same calculation at $w=0$.

The rational factors on the right of Equation~\eqref{eq:bulkboundarydeterminant} may have cancelling zeros and poles. The identity extends as the corresponding meromorphic identity; the product has the polynomial continuation on the left. Replacing the right side by $\det(\Id-wK_r)$ from the static boundary response generally discards both interior spectral data and the parameter dependence of $\mathsf S_r(w)$. A complete bulk determinant must retain these factors. A static boundary determinant is a legitimate \emph{different candidate}, but its equality to $F$ must be proved in its own right. Neither determinant is the Riemann function merely because it descends from a positive source quadratic form.

\subsection{Alternative arithmetic completion routes}
Theorems~\ref{thm:moment} and \ref{thm:primetrace} supply alternative conclusions of the construction. A positive trace-class realization of Equation~\eqref{eq:primebridge} already proves the determinant identity; no separate construction of $J_\sigma$ is then logically required for RH. Conversely, an all-degree positive representation of $q_\xi$ proves RH without first realizing the proposed boundary operator. Section~\ref{subsec:momentlimit} develops a finite-to-infinite version of the latter route. These alternatives retain the same missing requirement: exact agreement with the arithmetic target, rather than positivity of an unrelated model.

\section{Spectral growth required by the arithmetic target}
\label{sec:weyl}
\subsection{Counting the squared ordinates}
Suppose a positive determinant realization exists. Remove the kernel of $K$ and set $A=K^{-1}$ on its natural dense domain. Its eigenvalues are $\gamma_j^2$. The Riemann--von Mangoldt formula therefore requires
\begin{align}
 N_A(E)&=\#\{j:\gamma_j^2\leq E\}\notag\\
 &=\frac{\sqrt E}{4\pi}\log E
   -\frac{\log(2\pi)+1}{2\pi}\sqrt E+O(\log E).
 \label{eq:requiredcount}
\end{align}
This follows by substituting $T=\sqrt E$ into
$N(T)=\frac{T}{2\pi}\log\frac{T}{2\pi}-\frac{T}{2\pi}+O(\log T)$ \cite{titchmarsh}. Multiplicity is retained and only positive ordinates are counted. Doubling to include both signs would change every leading coefficient.

\begin{proposition}[Compact-elliptic obstruction]
\label{prop:weylobstruction}
The full spectrum of a positive classical elliptic operator of fixed order $m>0$ on a smooth compact $d$-dimensional manifold, under the usual regular elliptic boundary conditions when a boundary is present, cannot equal the squared zeta ordinates in a positive determinant realization. The same conclusion holds after a fixed positive affine spectral rescaling or a perturbation affecting only finitely many eigenvalues.
\end{proposition}
\begin{proof}
For this operator class the leading Weyl law is $N(E)\sim cE^{d/m}$, with $c>0$ determined by the principal symbol \cite{hormander}. Equation~\eqref{eq:requiredcount} instead has leading term $\frac1{4\pi}E^{1/2}\log E$. If $d/m\neq1/2$, the ratio of the two expressions tends to zero or infinity. If $d/m=1/2$, the logarithm still makes the ratio unbounded. An affine spectral rescaling changes coefficients but not this distinction, and finitely many eigenvalue changes alter $N(E)$ by $O(1)$.
\end{proof}

For a four-dimensional regular source chart, choosing order eight would match the power $E^{1/2}$ but not the logarithm. For a three-dimensional regular Shadow chart, order six has the same limitation. This is a restriction on an entire-spectrum realization in a specified operator class, not a dimensional refutation of Counterspace. Singular strata, noncompact corridors, nonlocal operators, logarithmic symbols, or independently selected infinite-dimensional subspaces fall outside the stated proposition. Each alternative requires its own counting theorem; the existence of singular source structure alone does not supply the required coefficient \cite{tcgsgeometry}.

\subsection{The heat-trace fingerprint}
The counting law determines an additional useful target. Here $\tau>0$ is an analytical heat parameter, not the operational clock of Equation~\eqref{eq:time}.

\begin{proposition}[Required heat asymptotics]
\label{prop:heat}
Any positive determinant realization has
\begin{equation}
 \Tr\e^{-\tau A}
 =\frac{\log(1/\tau)-\gamma_{\mathrm E}-2\log(4\pi)}
        {8\sqrt{\pi\tau}}
   +O\!\left(\log(1/\tau)\right)
 \quad(\tau\downarrow0),
 \label{eq:heat}
\end{equation}
where $\gamma_{\mathrm E}$ is Euler's constant.
\end{proposition}
\begin{proof}
Integration by parts gives
\[
 \Tr\e^{-\tau A}=\tau\int_0^\infty\e^{-\tau E}N_A(E)\dd E.
\]
Insert Equation~\eqref{eq:requiredcount}. With $v=\tau E$, the first two terms use $\Gamma(3/2)=\sqrt\pi/2$ and
$\Gamma'(3/2)/\Gamma(3/2)=2-\gamma_{\mathrm E}-2\log2$. Their coefficients simplify to Equation~\eqref{eq:heat}. Integrating the $O(\log E)$ remainder gives $O(\log(1/\tau))$; the bounded low-energy interval contributes only a bounded error.
\end{proof}

The leading logarithmic heat term is a stringent geometric requirement. It distinguishes the arithmetic target from an ordinary compact-manifold heat expansion even before individual eigenvalues are compared. It remains only a necessary condition: the exponential model below has the same first two counting terms and therefore the same displayed heat terms without having the correct determinant.

There is also a trace-ideal consequence. Inverting the leading zero-counting relation gives $\gamma_j\sim2\pi j/\log j$. Thus
\begin{equation}
 \lambda_j(K)\sim\left(\frac{\log j}{2\pi j}\right)^2,
 \qquad \Tr K^q<\infty\ \Longleftrightarrow\ q>\frac12.
 \label{eq:schatten}
\end{equation}
The borderline series at $q=1/2$ behaves like $\sum(\log j)/j$ and diverges. This calculation specifies how much compactness the arithmetic realization needs, rather than treating trace-class membership as an unspecified technical assumption.

\subsection{A smooth boundary response does not evade the obstruction}
\label{subsec:boundaryweyl}
Nonlocality by itself does not imply an anomalous leading counting law. For a smooth compact boundary, the usual Dirichlet-to-Neumann operator is a classical elliptic pseudodifferential operator of order one, and its spectrum is the Steklov spectrum \cite{steklov}. Its boundary dimension, operator order, and mass measure therefore remain decisive.

\begin{proposition}[Regular boundary-power obstruction]
\label{prop:boundarypower}
Let $B$ be such a positive boundary operator on a compact three-dimensional boundary, restricted off a finite-dimensional kernel when necessary, with eigenvalues $b_j\sim c j^{1/3}$, $c>0$. For any fixed $\alpha>0$, $B^{-\alpha}$ is trace class exactly when $\alpha>3$. No positive affine rescaling of $B^\alpha$ has the full counting law in Equation~\eqref{eq:requiredcount}.
\end{proposition}
\begin{proof}
The inverse-power series is comparable to $\sum_j j^{-\alpha/3}$, giving the trace-class condition. The counting function of $B^\alpha$ is asymptotic to a positive constant times $E^{3/\alpha}$. At $\alpha=6$ its power matches $E^{1/2}$ but the logarithm is absent; other powers fail already at the exponent. A positive affine rescaling does not change either conclusion.
\end{proof}

Thus an ordinary smooth boundary operator with a convenient inverse power cannot be promoted into the arithmetic realization. Singular interfaces, noncompact limits, or nonclassical spectral rules are possibilities to be analyzed, not guaranteed solutions. The compact-elliptic and calibrated Morse exclusions remain restricted to their stated classes; they do not exclude every smooth potential or all conventional spectral mathematics. In particular, the noncompact smooth exponential model below already reproduces the logarithmic density while failing the exact determinant.

\section{An explicit exponential-corridor construction}
\label{sec:exponential}
\subsection{The positive operator}
Consider the differential expression
\begin{equation}
 A_{\mathrm{exp}}=-\frac{\dd^2}{\dd u^2}+16\pi^2\e^{4u}
 \quad\text{on }L^2([0,\infty),\dd u),\qquad f(0)=0.
 \label{eq:expoperator}
\end{equation}
Its coefficient is calibrated to the required asymptotic density, not derived from the ECL. The variable $u$ is an auxiliary logarithmic corridor coordinate. No physical one-dimensional source or ontic time is asserted. This example tests whether a particularly economical noncompact realization can satisfy the arithmetic conditions left open by the compact-elliptic obstruction.

Define the closed quadratic form
\begin{equation}
 a_{\mathrm{exp}}[f]=\int_0^\infty
 \left(|f'(u)|^2+16\pi^2\e^{4u}|f(u)|^2\right)\dd u,
 \label{eq:expform}
\end{equation}
with form domain $H_0^1([0,\infty))\cap L^2([0,\infty),\e^{4u}\dd u)$. Its associated Friedrichs operator is the realization intended in Equation~\eqref{eq:expoperator}.

\begin{proposition}[Well-defined positive corridor]
\label{prop:exppositive}
$A_{\mathrm{exp}}$ is self-adjoint, satisfies $A_{\mathrm{exp}}\geq16\pi^2\Id$, and has compact resolvent. Its eigenvalues are positive, discrete, and simple.
\end{proposition}
\begin{proof}
The densely defined closed form is bounded below by $16\pi^2\|f\|_2^2$, so the representation theorem for semibounded forms gives a self-adjoint operator with that bound \cite{reed}. To prove compactness, consider a set bounded in the form norm. On each interval $[0,R]$, the $H^1$ bound gives compactness in $L^2$. The tail satisfies
\[
 \int_R^\infty|f(u)|^2\dd u
 \leq\frac{\e^{-4R}}{16\pi^2}a_{\mathrm{exp}}[f],
\]
so it is uniformly small. A diagonal subsequence argument gives compactness of the full form-domain embedding, hence compact resolvent. Simplicity follows from the Wronskian: two square-integrable solutions at infinity with the same eigenvalue are proportional, and the Dirichlet condition supplies only one independent endpoint condition.
\end{proof}

\subsection{An exactly computable counting law}
For $E>16\pi^2$, let $\delta=4\pi/\sqrt E$ and let $u_E=\frac14\log(E/(16\pi^2))$ be the turning point. The classical one-dimensional phase integral is
\begin{align}
 I_{\mathrm{exp}}(E)
 &=\frac1\pi\int_0^{u_E}\sqrt{E-16\pi^2\e^{4u}}\dd u\notag\\
 &=\frac{\sqrt E}{2\pi}
 \left[\log\frac{1+\sqrt{1-\delta^2}}{\delta}
       -\sqrt{1-\delta^2}\right].
 \label{eq:phaseintegral}
\end{align}
The second equality follows from the substitution $v=4\pi\e^{2u}/\sqrt E$ and integration of $\sqrt{1-v^2}/v$ between $\delta$ and one.

For this smooth monotone exponential potential, the one-dimensional counting theorem gives $N_{\mathrm{exp}}(E)=I_{\mathrm{exp}}(E)+O(1)$. It is also a direct rescaling of the Morse half-line counting result: setting $r=2u+\log(4\pi)$ gives four times the operator $-\dd^2/\dd r^2+\frac14\e^{2r}$ on $r\geq\log(4\pi)$ \cite{lagarias}. Expanding Equation~\eqref{eq:phaseintegral} therefore gives
\begin{equation}
 N_{\mathrm{exp}}(E)
 =\frac{\sqrt E}{2\pi}
   \left(\log\frac{\sqrt E}{2\pi}-1\right)+O(1).
 \label{eq:expcount}
\end{equation}
Both displayed coefficients agree with Equation~\eqref{eq:requiredcount}. The available classically allowed corridor length grows logarithmically with $E$, which is the explicit mechanism generating the extra logarithm absent from the compact Weyl law.

This count implies $\sum_jE_j^{-1}<\infty$, for example by integrating $E^{-1}\dd N_{\mathrm{exp}}(E)$ at infinity. Thus $A_{\mathrm{exp}}^{-1}$ is positive trace class. The construction has now met positivity, discreteness, trace-class compactness, and the first two arithmetic counting terms. It remains to test the entire determinant.

\subsection{The Bessel characteristic determinant}
Put $y=2\pi\e^{2u}$. The eigenvalue equation reduces to
\begin{equation}
 y^2f_{yy}+yf_y+\left(\frac E4-y^2\right)f=0.
\end{equation}
The solution decaying at infinity is proportional to the modified Bessel function $K_{i\sqrt E/2}(y)$. Consequently the Dirichlet characteristic function is
\begin{equation}
 D_{\mathrm{exp}}(w)
 =\frac{K_{i\sqrt w/2}(2\pi)}{K_0(2\pi)}.
 \label{eq:besseldet}
\end{equation}
The evenness of $K_\nu$ in its order makes this an entire function of $w$. Bessel-order zeros and their spectral interpretation have classical antecedents; this is a concrete normalization of that family \cite{lagarias,dlmf}.

\begin{proposition}[Exact determinant of the corridor]
\label{prop:expdet}
The characteristic function satisfies
\begin{equation}
 D_{\mathrm{exp}}(w)=\det(\Id-wA_{\mathrm{exp}}^{-1}).
 \label{eq:expfredholm}
\end{equation}
\end{proposition}
\begin{proof}
A zero of the numerator in Equation~\eqref{eq:besseldet} supplies a decaying Dirichlet solution and is therefore an eigenvalue. Conversely each eigenvalue gives such a zero. The zeros are simple: differentiating the differential equation in $w$ and applying Green's identity at an eigenvalue $E$ gives
\[
 \int_0^\infty f(u,E)^2\dd u
 =-f_u(0,E)\,\partial_w f(0,E).
\]
The left side is positive, and $f_u(0,E)\neq0$ by uniqueness for the second-order equation. Thus the boundary characteristic has nonzero derivative.

The integral $K_{i\sqrt w/2}(2\pi)=\int_0^\infty\e^{-2\pi\cosh v}\cos(\sqrt w\,v/2)\dd v$ bounds its maximum modulus on $|w|\leq R$ by $K_{\sqrt R/2}(2\pi)$. The large-order bound is $\log K_{\sqrt R/2}(2\pi)=O(\sqrt R\log R)$, so the entire order is at most $1/2$ \cite{dlmf}. The counting law forces order $1/2$. Its normalized genus-zero product is therefore exactly the product over eigenvalues in Equation~\eqref{eq:expfredholm}.
\end{proof}

\subsection{The determinant fails despite the correct density}
\begin{theorem}[Exponential determinant mismatch]
\label{thm:expmismatch}
For real $s\to+\infty$,
\begin{equation}
 \frac{F(-(s-\tfrac12)^2)}{D_{\mathrm{exp}}(-(s-\tfrac12)^2)}
 \sim\frac{K_0(2\pi)}{\xi(\tfrac12)(2\pi)^{1/4}}s^{9/4}.
 \label{eq:expmismatch}
\end{equation}
In particular $D_{\mathrm{exp}}\neq F$, and $A_{\mathrm{exp}}^{-1}$ does not satisfy the prime-trace identity.
\end{theorem}
\begin{proof}
Write $\nu=s/2-1/4$. At fixed positive argument, the Bessel integral can be written
\[
 K_\nu(2\pi)=\frac12\pi^{-\nu}
 \int_0^\infty y^{\nu-1}\e^{-y-\pi^2/y}\dd y.
\]
For a gamma random variable $Y$ of shape $\nu>1$ and unit scale, this gives
\[
 \frac{K_\nu(2\pi)}{\frac12\pi^{-\nu}\Gamma(\nu)}
 =\mathbb E\e^{-\pi^2/Y}.
\]
Since $0\leq1-\e^{-v}\leq v$,
\begin{equation}
 1-\frac{\pi^2}{\nu-1}
 \leq\mathbb E\e^{-\pi^2/Y}\leq1.
 \label{eq:besselbound}
\end{equation}
Thus $K_\nu(2\pi)\sim\frac12\pi^{-\nu}\Gamma(\nu)$; the same standard asymptotic is recorded in the Bessel literature \cite{dlmf}.

For real $s\to\infty$, $\zeta(s)=1+O(2^{-s})$, so Equation~\eqref{eq:xi} gives
\[
 \xi(s)\sim\frac12s(s-1)\pi^{-s/2}\Gamma(s/2).
\]
Use $F(-(s-1/2)^2)=\xi(s)/\xi(1/2)$, the evenness of Bessel order, and
$\Gamma(s/2)/\Gamma(s/2-1/4)\sim(s/2)^{1/4}$. The quotient simplifies to Equation~\eqref{eq:expmismatch}. An identity of determinants would make the quotient identically one, which is impossible. The final assertion follows from Theorem~\ref{thm:primetrace}.
\end{proof}

This failure is stronger than a numerical disagreement in a low eigenvalue. It is an analytical obstruction in an asymptotic regime where all quantities are real and positive. The model does not merely need a better estimate of its existing spectrum: its characteristic function has the wrong normalization-dependent power. The next construction tests whether a minimal additional potential term can repair this defect.

\section{A Morse correction and a family-level obstruction}
\label{sec:morse}
\subsection{The one-parameter repair}
Consider the calibrated Dirichlet family
\begin{equation}
 A_\kappa=-\frac{\dd^2}{\dd u^2}
     +16\pi^2\e^{4u}-16\pi\kappa\e^{2u},
 \qquad u\geq0,\quad f(0)=0,\quad\kappa<\pi.
 \label{eq:morseoperator}
\end{equation}
Here $\kappa$ is a model parameter, not a Counterspace invariant. For $y=\e^{2u}\geq1$, the potential is $16\pi y(\pi y-\kappa)$; it is increasing and bounded below by $16\pi(\pi-\kappa)>0$. The form argument used for Proposition~\ref{prop:exppositive} therefore gives a positive self-adjoint operator with compact resolvent.

Under $r=2u+\log(4\pi)$, Equation~\eqref{eq:morseoperator} becomes four times the standard Morse operator with potential $\frac14\e^{2r}-\kappa\e^r$.\footnote{Lagarias writes $V_k(r)=\tfrac14\e^{2r}+k\e^r$, so his parameter is $k=-\kappa$ in this normalized expression, and the decaying solution uses $W_{-k,\nu}=W_{\kappa,\nu}$. The spectral factor four and the change of endpoint must also be retained when applying his counting theorem \cite[Theorems~2.1 and 3.2]{lagarias}.} Its eigenvalue count retains Equation~\eqref{eq:expcount}, with a $\kappa$-dependent bounded remainder, and its inverse is trace class \cite{lagarias}. The decaying solution is proportional to
\begin{equation}
 (4\pi\e^{2u})^{-1/2}
 W_{\kappa,i\sqrt w/2}(4\pi\e^{2u}),
\end{equation}
where $W$ is the Whittaker function. The normalized determinant is
\begin{equation}
 D_\kappa(w)=\frac{W_{\kappa,i\sqrt w/2}(4\pi)}{W_{\kappa,0}(4\pi)}
           =\det(\Id-wA_\kappa^{-1}).
 \label{eq:morsedet}
\end{equation}
The denominator is nonzero because zero is not an eigenvalue. Lemma~\ref{lem:complexwhittakergrowth} below supplies an explicit bound throughout the complex order plane, and Proposition~\ref{prop:whittakerproducts} proves the simple-zero and normalized genus-zero product statements, for this Dirichlet case and for fixed Robin data. The real-order asymptotic of Lemma~\ref{lem:whittaker} is used only for the subsequent comparison, not as a complex-plane growth bound.

\subsection{The necessary asymptotic correction}
The following coefficient, derived in Appendix~\ref{app:whittaker}, is sufficient to test every member of the family.

\begin{lemma}[Fixed-argument Whittaker expansion]
\label{lem:whittaker}
For fixed $x>0$ and fixed real $\kappa$, as $\mu\to+\infty$,
\begin{align}
 W_{\kappa,\mu}(x)
 &=\frac{\sqrt x}{2\sqrt\pi}\Gamma(\mu)
       (x/4)^{-\mu}\mu^\kappa
       \left[1+\frac{c(\kappa,x)}{\mu}+O(\mu^{-2})\right],\label{eq:whittakerexpansion}\\
 c(\kappa,x)&=-\frac{\kappa^2}{2}+\frac{\kappa x}{2}-\frac{x^2}{16}.
 \label{eq:whittakercoefficient}
\end{align}
\end{lemma}

\begin{theorem}[Calibrated Dirichlet Morse-family exclusion]
\label{thm:morseexclusion}
No member of Equation~\eqref{eq:morseoperator}, for $\kappa<\pi$, has $D_\kappa=F$. More precisely, with $\mu=s/2-1/4$,
\begin{equation}
 \frac{F(-(s-\tfrac12)^2)}{D_\kappa(-(s-\tfrac12)^2)}
 =C_\kappa\mu^{9/4-\kappa}
 \left[1+\frac{\delta_\kappa}{\mu}+O(\mu^{-2})\right],
 \label{eq:morseratio}
\end{equation}
where
\begin{equation}
 C_\kappa=\frac{2W_{\kappa,0}(4\pi)}{\xi(\tfrac12)\pi^{1/4}}>0,
 \qquad
 \delta_\kappa=\frac{\kappa^2}{2}-2\pi\kappa+\pi^2-\frac3{32}.
 \label{eq:morseconstants}
\end{equation}
The power mismatch can vanish only at $\kappa=9/4$, but then
\begin{equation}
 \delta_{9/4}=\pi^2-\frac92\pi+\frac{39}{16}<0,
 \label{eq:nonzeroresidue}
\end{equation}
which prevents the quotient from being identically one.
\end{theorem}
\begin{proof}
At $s=2\mu+1/2$, the defining completion gives
\[
 \xi(s)=\left(2\mu^2-\frac18\right)
 \pi^{-\mu-1/4}\Gamma(\mu+1/4)\zeta(2\mu+1/2).
\]
The gamma-ratio expansion and the exponentially small zeta correction yield
\begin{equation}
 \xi(2\mu+\tfrac12)
 =2\pi^{-1/4}\Gamma(\mu)\pi^{-\mu}\mu^{9/4}
   \left[1-\frac3{32\mu}+O(\mu^{-2})\right].
 \label{eq:xiorderexpansion}
\end{equation}
At $x=4\pi$, Lemma~\ref{lem:whittaker} has leading prefactor $\Gamma(\mu)\pi^{-\mu}\mu^\kappa$ and correction
$-\kappa^2/2+2\pi\kappa-\pi^2$. Dividing Equation~\eqref{eq:xiorderexpansion} by that expansion, with the normalizations in Equation~\eqref{eq:morsedet}, proves Equations~\eqref{eq:morseratio}--\eqref{eq:morseconstants}.

For $\kappa\neq9/4$, the quotient tends to zero or infinity. For $\kappa=9/4$, an identical quotient would first require $C_{9/4}=1$ and then $\delta_{9/4}=0$. Equation~\eqref{eq:nonzeroresidue} rules out the second requirement, independently of the value of the first. Its sign follows, for example, from $3<\pi<22/7$. Finally, $W_{\kappa,0}(4\pi)>0$: its decaying real solution is positive at infinity and cannot have a zero on the half-line, since such a zero would produce a zero-energy Dirichlet eigenfunction on a shorter interval with strictly positive potential. Hence $C_\kappa>0$.
\end{proof}

The family-level conclusion is deliberately bounded. It excludes the calibrated exponential-plus-one-Morse-term potential with the stated Dirichlet endpoint, not every singular Sturm--Liouville system, every boundary condition, or every source-derived operator. Its strength is that it tests a natural attempted repair rather than stopping at the first mismatch. The correction $\kappa=9/4$ removes the leading power defect but cannot reproduce the next coefficient of the actual arithmetic completion.

At the framework level, this provides a precise example of the distinction between a transferable constitutive map and an expressive fitting family. The parameter was introduced to repair a known target asymptotic. Even a successful repair would not have been a source derivation without an independent rule selecting it. The stronger result obtained here is that this repair fails the arithmetic identity on its own mathematical terms. Neither source completeness nor the positivity of the potential changes the nonzero coefficient in Equation~\eqref{eq:nonzeroresidue}.


\subsection{Free potential scale and endpoint}
\label{subsec:scaleendpoint}
The Dirichlet endpoint and exponential scale need not be fixed independently. With the kinetic coefficient and exponential rate retained as in Equation~\eqref{eq:morseoperator}, consider
\begin{equation}
 A_{c,\kappa,u_0}=-\frac{\dd^2}{\dd u^2}
       +c^2\e^{4u}-4c\kappa\e^{2u},\qquad
 u\in[u_0,\infty),\quad f(u_0)=0,
 \quad c>0,\quad \kappa,u_0\in\R.
 \label{eq:freemorse}
\end{equation}
The change of coordinate $r=2u+\log c$, including the unitary factor $2^{-1/2}$ on the wave function, gives
\begin{equation}
 A_{c,\kappa,u_0}\simeq
 4\left(-\frac{\dd^2}{\dd r^2}
             +\frac14\e^{2r}-\kappa\e^r\right),\qquad
 r\geq\log x,\qquad x:=c\e^{2u_0}>0.
 \label{eq:effectiveendpoint}
\end{equation}
Thus scale and endpoint are not two independent spectral adjustments: for fixed $\kappa$, their only spectral combination is $x$. The classical Morse counting law \cite[Theorem~3.2]{lagarias}, with the factor four in Equation~\eqref{eq:effectiveendpoint}, reads
\begin{equation}
 N_{c,\kappa,u_0}(E)
 =\frac{\sqrt E}{2\pi}
   \left[\log\sqrt E+\log\frac2x-1\right]+O(1).
 \label{eq:freeendpointcount}
\end{equation}
Comparison with Equation~\eqref{eq:requiredcount} forces
\begin{equation}
 \log\frac2x=-\log(2\pi),\qquad
 x=c\e^{2u_0}=4\pi.
 \label{eq:endpointcalibration}
\end{equation}
The $O(1)$ correction depends on the fixed parameters but cannot change the coefficient of $\sqrt E$.

\begin{corollary}[Scale--endpoint closure of the Dirichlet Morse exclusion]
\label{cor:scaleendpointexclusion}
No strictly positive Dirichlet realization of Equation~\eqref{eq:freemorse} has normalized determinant $F$. Hence freeing the potential scale and endpoint does not repair the exclusion in Theorem~\ref{thm:morseexclusion}.
\end{corollary}
\begin{proof}
Every operator in Equation~\eqref{eq:freemorse} is semibounded with compact resolvent, since its potential is $(c\e^{2u}-2\kappa)^2-4\kappa^2$ and tends to positive infinity. In the strictly positive case, Equation~\eqref{eq:freeendpointcount} implies that the inverse is trace class. The same characteristic-function argument as in Equation~\eqref{eq:morsedet} gives
\begin{equation}
 D_{\kappa,x}(w)
 =\frac{W_{\kappa,i\sqrt w/2}(x)}{W_{\kappa,0}(x)}
 =\det(\Id-wA_{c,\kappa,u_0}^{-1}).
 \label{eq:freeendpointdet}
\end{equation}
The denominator is nonzero. Put $\mu=s/2-1/4$ for real $s\to+\infty$. Lemma~\ref{lem:whittaker} and Equation~\eqref{eq:xiorderexpansion} yield
\begin{align}
 \frac{F(-(s-\tfrac12)^2)}{D_{\kappa,x}(-(s-\tfrac12)^2)}
 &=C_{\kappa,x}
   \left(\frac{x}{4\pi}\right)^\mu
   \mu^{9/4-\kappa}
   \left[1+\frac{\delta_\kappa(x)}{\mu}+O(\mu^{-2})\right],
   \label{eq:freeendpointratio}\\
 C_{\kappa,x}&=\frac{4\pi^{1/4}W_{\kappa,0}(x)}{\sqrt x\,\xi(\tfrac12)},
   \qquad
 \delta_\kappa(x)=\frac{\kappa^2}{2}-\frac{\kappa x}{2}
                    +\frac{x^2}{16}-\frac3{32}.
   \label{eq:freeendpointconstants}
\end{align}
The constant is positive in the strictly positive realization, as also follows by comparing the positive determinants on the negative real axis. An identity of determinants first forces $x=4\pi$ by the exponential factor, consistently with the independent counting argument. It then forces $\kappa=9/4$ by the power factor. At that value,
\begin{equation}
 \delta_{9/4}(x)=\frac{x^2-18x+39}{16},\qquad
 \delta_{9/4}(x)=0\ \Longleftrightarrow\ x=9\pm\sqrt{42}.
 \label{eq:endpointroots}
\end{equation}
Neither root equals $4\pi$, where Equation~\eqref{eq:nonzeroresidue} gives the nonzero coefficient already computed. Thus equality is impossible, regardless of the leading normalization constant.
\end{proof}

The two roots are approximately $2.5192593$ and $15.4807407$. Canceling the inverse-order coefficient at either root breaks the necessary second counting coefficient and leaves the exponential factor in Equation~\eqref{eq:freeendpointratio}. At fixed $x=4\pi$, changing $c$ and compensating with $u_0$ merely gives a unitarily equivalent operator. The exclusion concerns this full scale--endpoint Dirichlet family with the stated spectral normalization, not arbitrary boundary conditions or additional potentials. A realization with a zero or negative eigenvalue cannot supply a positive inverse determinant with zero set $(0,\infty)$; deleting such modes or adding a spectral shift defines a different candidate and is not covered by Corollary~\ref{cor:scaleendpointexclusion} alone. Constant spectral shifts and all fixed real separated endpoint conditions are addressed explicitly in Sections~\ref{subsec:constantshift} and \ref{subsec:robinexclusion}; deletion of modes remains outside those statements.


\subsection{Complex-order growth and canonical determinant normalization}
\label{subsec:complexwhittaker}
A real-order asymptotic does not control an entire function in the complex spectral plane. The determinant identification in Equations~\eqref{eq:morsedet} and \eqref{eq:freeendpointdet} requires the latter control. The following estimate displays it and also covers a fixed Robin endpoint. The special-function inputs are the Tricomi integral, the Whittaker--Tricomi relation, evenness in the order, and the reciprocal-gamma product \cite[Sections~5.8, 13.4, 13.14, and 13.16]{dlmf}; no zero-location assumption about $\xi$ is used.

\begin{lemma}[Uniform complex-order growth]
\label{lem:complexwhittakergrowth}
Fix $\kappa\in\R$ and $0<x_-<x_+<\infty$. There is a constant $C$ such that, for $R\geq1$,
\begin{equation}
 \sup_{\substack{|\nu|\leq R\\x\in[x_-,x_+]}}
 \left(|W_{\kappa,\nu}(x)|+|\partial_xW_{\kappa,\nu}(x)|\right)
 \leq \exp\!\left(C(1+R)\log(2+R)\right).
 \label{eq:complexwhittakerbound}
\end{equation}
Consequently $W_{\kappa,i\sqrt w/2}(x)$ and
$2x\partial_xW_{\kappa,i\sqrt w/2}(x)-(1+h)W_{\kappa,i\sqrt w/2}(x)$, for fixed real $h$, are entire functions of $w$ of order at most $1/2$. The square-root notation denotes their even power series.
\end{lemma}
\begin{proof}
Choose an integer $m\geq0$ such that $b_0=m+1/2-\kappa>0$, and put
$a=\nu+1/2-\kappa$ and $b=\nu+\kappa-1/2$. For $\Re a>0$, the standard integral and $m$ integrations by parts give
\begin{align}
 W_{\kappa,\nu}(x)
 &=\frac{\e^{-x/2}x^{\nu+1/2}}{\Gamma(a)}
   \int_0^\infty t^{a-1}\e^{-xt}(1+t)^b\dd t\notag\\
 &=\frac{(-1)^m\e^{-x/2}x^{\nu+1/2}}{\Gamma(a+m)}
   \int_0^\infty t^{a+m-1}
   \frac{\dd^m}{\dd t^m}\!\left[\e^{-xt}(1+t)^b\right]\dd t.
 \label{eq:continuedwhittakerintegral}
\end{align}
The boundary terms vanish in the initial half-plane. The last integral is holomorphic whenever $\Re\nu>-b_0$, so analytic continuation makes it valid throughout $\Re\nu\geq0$, including cases where the first integral is not convergent. Its differentiated factor is the finite sum
\begin{equation}
 \frac{\dd^m}{\dd t^m}\!\left[\e^{-xt}(1+t)^b\right]
 =\e^{-xt}\sum_{j=0}^{m}\binom mj(-x)^{m-j}
       b^{\underline j}(1+t)^{b-j},
 \qquad b^{\underline j}=\prod_{\ell=0}^{j-1}(b-\ell).
 \label{eq:whittakerfiniteparts}
\end{equation}
For $p=\Re\nu\in[0,R]$, its coefficients have magnitude at most $C_1(1+R)^m$. Splitting the integral at $t=1$ bounds the remaining absolute integral, uniformly in $x$, by
\begin{equation}
 C_2\,2^R\left(1+x_-^{-2R-B}\Gamma(2R+B)\right)
 \label{eq:whittakerintegralmajorant}
\end{equation}
for a fixed sufficiently large $B>0$. On $(0,1)$ one uses $t^{p+b_0-1}\leq t^{b_0-1}$; on $[1,\infty)$ the powers are bounded by a fixed multiple of $2^Rt^{2R+B-1}$. The logarithm of the right side is $O((1+R)\log(2+R))$.

The remaining reciprocal gamma factor has the same required bound. Indeed,
\begin{equation}
 \Gamma(z)^{-1}=z\e^{\gamma_{\mathrm E}z}
      \prod_{n=1}^{\infty}(1+z/n)\e^{-z/n}.
 \label{eq:reciprocalgammaproduct}
\end{equation}
For $|z|\leq R+C_3$, split the product at an integer exceeding $2(R+C_3)$. The logarithm of the modulus of the finite part is bounded above by
$\sum[\log(1+|z|/n)+|z|/n]=O((1+R)\log(2+R))$.
The tail is bounded by $C_4|z|^2\sum n^{-2}=O(1+R)$, using $|z/n|\leq1/2$. Zeros of the reciprocal gamma function cause no difficulty for this upper bound. Also $|x^{\nu+1/2}|\leq C_5\exp(C_6R)$.

Differentiating Equation~\eqref{eq:continuedwhittakerintegral} in $x$ only adds a factor bounded by $C_7(1+R)$ outside the integral and finitely many terms with an extra factor $t$ inside it. These are bounded by increasing $B$ in Equation~\eqref{eq:whittakerintegralmajorant}. Differentiation under the integral is justified by the same exponential majorant on the compact $x$-interval. This proves Equation~\eqref{eq:complexwhittakerbound} in the right half-plane. The standard evenness $W_{\kappa,-\nu}=W_{\kappa,\nu}$, also after $x$-differentiation, supplies the other half-plane. Entire dependence on $\nu$ and this evenness give entire dependence on $w$; substituting $|\nu|\leq\sqrt{|w|}/2$ gives order at most $1/2$.
\end{proof}

\begin{proposition}[Dirichlet and Robin characteristic products, including shifts]
\label{prop:whittakerproducts}
Let $A$ be the realization of
\begin{equation}
 -\frac{\dd^2}{\dd u^2}+c^2\e^{4u}-4c\kappa\e^{2u}+\lambda,
 \qquad u\geq u_0,
 \label{eq:shiftedmorsedifferential}
\end{equation}
with $c>0$, real fixed $\kappa,u_0,\lambda$, and either $f(u_0)=0$ or $f'(u_0)=h f(u_0)$ for a fixed $h\in\R$. Assume $A$ is strictly positive, and set $x=c\e^{2u_0}$. Define the unshifted characteristic functions
\begin{align}
 \mathcal H_D(w)&=W_{\kappa,i\sqrt w/2}(x),\notag\\
 \mathcal H_R(w)&=Q_{\kappa,h,i\sqrt w/2}(x),\qquad
 Q_{\kappa,h,\nu}(x)=2x\partial_xW_{\kappa,\nu}(x)-(1+h)W_{\kappa,\nu}(x).
 \label{eq:robincharacteristic}
\end{align}
For the chosen boundary condition $B\in\{D,R\}$, one has
\begin{equation}
 \frac{\mathcal H_B(w-\lambda)}{\mathcal H_B(-\lambda)}
 =\prod_{j=1}^{\infty}(1-w/E_j)
 =\det(\Id-wA^{-1}),
 \label{eq:shiftedcharacteristicproduct}
\end{equation}
where $E_j>0$ are the eigenvalues of $A$, without deletion. The product has genus zero and the entire function has order $1/2$.
\end{proposition}
\begin{proof}
The real separated half-line Morse problem is self-adjoint, has compact resolvent, and has simple discrete spectrum. For every fixed endpoint condition its counting law is Equation~\eqref{eq:freeendpointcount} \cite[Theorems~2.1 and 3.2]{lagarias}. A constant spectral shift leaves its two displayed coefficients unchanged, so $\sum E_j^{-1}<\infty$ under the stated positivity assumption.

With $y=c\e^{2u}$, the decaying solution is
$f(u,w)=y^{-1/2}W_{\kappa,\nu}(y)$ with $\nu^2=(\lambda-w)/4$.
Its Dirichlet boundary value is $x^{-1/2}\mathcal H_D(w-\lambda)$; its Robin residual is
$x^{-1/2}\mathcal H_R(w-\lambda)$. Thus each characteristic zero is an eigenvalue and conversely. To check analytic multiplicity, take the real decaying solution at a real eigenvalue $E$ and put $g=\partial_w f$. Let $\mathcal L$ denote the differential expression in Equation~\eqref{eq:shiftedmorsedifferential}. Differentiating before imposing any endpoint condition on $g$ gives $(\mathcal L-E)g=f$. Green's identity gives
\begin{equation}
 \int_{u_0}^{\infty} f(u,E)^2\dd u
 = f(u_0,E)g'(u_0,E)-f'(u_0,E)g(u_0,E).
 \label{eq:robinmultiplicity}
\end{equation}
The terms at infinity vanish by the decaying Whittaker asymptotic and its parameter derivative. For Dirichlet data this equals $-f'(u_0,E)g(u_0,E)$; for Robin data it equals $f(u_0,E)(g'(u_0,E)-hg(u_0,E))$. The nonzero endpoint factor in either case cannot vanish by uniqueness for the differential equation. The positive integral therefore proves that the characteristic zero is simple.

Lemma~\ref{lem:complexwhittakergrowth} bounds the full complex-plane growth by
$\exp(O(\sqrt{|w|}\log(2+|w|)))$; a fixed translation $w\mapsto w-\lambda$ preserves this order bound. Hadamard factorization for order less than one supplies no nonconstant exponential factor. Since zero is not an eigenvalue, division by $\mathcal H_B(-\lambda)$ fixes the remaining constant to one. The trace-class product is exactly Equation~\eqref{eq:shiftedcharacteristicproduct}. The eigenvalue count gives the matching lower bound on entire order.
\end{proof}

\subsection{Constant spectral shifts}
\label{subsec:constantshift}
A fixed spectral shift is invisible to the first two counting coefficients but not to the full determinant. More precisely, the change of the \emph{smooth counting main term} under $E\mapsto E-\lambda$ is $O(E^{-1/2}\log E)$. This is not a vanishing bound on the difference of two integer-valued counting staircases; the unchanged $O(1)$ counting remainder must still be retained.

\begin{corollary}[Constant-shift Dirichlet exclusion]
\label{cor:shiftedmorseexclusion}
No strictly positive Dirichlet realization of Equation~\eqref{eq:shiftedmorsedifferential}, for any real constant $\lambda$, has normalized inverse determinant $F$.
\end{corollary}
\begin{proof}
Write $s=2\mu+1/2$. At $w=-4\mu^2$, the shifted Whittaker order is
\begin{equation}
 \mu_\lambda=\sqrt{\mu^2+\lambda/4}
 =\mu+\frac{\lambda}{8\mu}+O(\mu^{-3}).
 \label{eq:shiftedorder}
\end{equation}
The gamma-ratio expansion, or the mean-value formula with
$\psi_\Gamma(\mu)=\log\mu+O(\mu^{-1})$, gives
\begin{equation}
 \log\frac{\Gamma(\mu_\lambda)(x/4)^{-\mu_\lambda}
                         \mu_\lambda^\kappa}
                    {\Gamma(\mu)(x/4)^{-\mu}\mu^\kappa}
 =\frac{\lambda}{8\mu}\bigl(\log\mu-\log(x/4)\bigr)
   +O(\mu^{-2}).
 \label{eq:shiftlogterm}
\end{equation}
Apply Lemma~\ref{lem:whittaker} and Proposition~\ref{prop:whittakerproducts}. Equality to $F$ first forces $x=4\pi$ and $\kappa=9/4$ by the exponential and power factors, respectively, and fixes the leading normalization. The remaining logarithm of the quotient is
\begin{equation}
 -\frac{\lambda}{8\mu}(\log\mu-\log\pi)
       +\frac{\delta_{9/4}(4\pi)}{\mu}+O(\mu^{-2}).
 \label{eq:shiftedmorsecontradiction}
\end{equation}
Multiplication by $\mu/\log\mu$ forces $\lambda=0$. The nonzero coefficient in Equation~\eqref{eq:nonzeroresidue} then contradicts equality. The initial normalization is $\mathcal H_D(-\lambda)$, not $\mathcal H_D(0)$, so the argument also covers shifts of unshifted operators having a zero or negative eigenvalue.
\end{proof}

\subsection{Robin matching and an analytic second-order obstruction}
\label{subsec:robinexclusion}
Robin data are not covered by the Dirichlet power comparison. The derivative in Equation~\eqref{eq:robincharacteristic} supplies an additional power of the large order. To track it beyond leading order, define
\begin{equation}
 L_{\kappa,\mu}(x)=\frac{\sqrt x}{2\sqrt\pi}
                    \Gamma(\mu)(x/4)^{-\mu}\mu^\kappa.
 \label{eq:whittakerleadingfactor}
\end{equation}
Appendix~\ref{subsec:secondorderwhittaker} proves the locally $x$-uniform expansion, including its first $x$-derivative,
\begin{align}
 \log\frac{W_{\kappa,\mu}(x)}{L_{\kappa,\mu}(x)}
 &=\frac{c(\kappa,x)}{\mu}+\frac{c_2(\kappa,x)}{\mu^2}
   +O(\mu^{-3}),\notag\\
 c_2(\kappa,x)&=-\frac{\kappa^3}{6}+\frac{\kappa}{24}
                  +\frac{\kappa x}{4}-\frac{x^2}{16},
 \label{eq:whittakerlogsecond}
\end{align}
with $c(\kappa,x)$ as in Equation~\eqref{eq:whittakercoefficient}. In particular,
\begin{equation}
 \frac{\partial_xW_{\kappa,\mu}(x)}{W_{\kappa,\mu}(x)}
 =-\frac\mu x+\frac1{2x}
   +\frac{\kappa/2-x/8}{\mu}
   +\frac{\kappa/4-x/8}{\mu^2}+O(\mu^{-3}).
 \label{eq:whittakerlogderivative}
\end{equation}
This gives
\begin{equation}
 \frac{Q_{\kappa,h,\mu}(x)}{-2\mu W_{\kappa,\mu}(x)}
 =1+\frac{h}{2\mu}
     -\frac{\kappa x/2-x^2/8}{\mu^2}+O(\mu^{-3}).
 \label{eq:robinprefactor}
\end{equation}
The target, in the same order variable, satisfies
\begin{equation}
 \log\frac{\xi(2\mu+1/2)}
 {2\pi^{-1/4}\Gamma(\mu)\pi^{-\mu}\mu^{9/4}}
 =-\frac3{32\mu}-\frac9{128\mu^2}+O(\mu^{-3}).
 \label{eq:xilogsecond}
\end{equation}
Thus the finite Robin parameter can cancel the first inverse-order mismatch after the power has been matched. It cannot be omitted from the comparison.

\begin{theorem}[Constant Robin and spectral-shift exclusion]
\label{thm:robinmorseexclusion}
No strictly positive realization of Equation~\eqref{eq:shiftedmorsedifferential} with real constant Robin condition $f'(u_0)=hf(u_0)$ has normalized inverse determinant $F$. This holds for every $c>0$ and fixed real $\kappa,u_0,h,\lambda$, with the kinetic coefficient and exponential rate displayed in that equation.
\end{theorem}
\begin{proof}
For $\lambda=0$, write $D_R(w)=\mathcal H_R(w)/\mathcal H_R(0)$ in the positive realization. Equations~\eqref{eq:whittakerlogsecond}--\eqref{eq:xilogsecond} give the full real-axis comparison
\begin{align}
 \log\frac{F(-4\mu^2)}{D_R(-4\mu^2)}
 &=\log C_R+\mu\log\frac{x}{4\pi}
       +(5/4-\kappa)\log\mu\notag\\
 &\quad+\frac{\delta_\kappa(x)-h/2}{\mu}
       +\frac{b_{\kappa,h}(x)}{\mu^2}+O(\mu^{-3}),
 \label{eq:robinlogratio}\\
 C_R&=-\frac{2\pi^{1/4}Q_{\kappa,h,0}(x)}{\sqrt x\,\xi(1/2)}>0,\notag\\
 b_{\kappa,h}(x)&=-\frac9{128}+\frac{\kappa^3}{6}-\frac{\kappa}{24}
                  +\frac{\kappa x}{4}-\frac{x^2}{16}+\frac{h^2}{8}.
 \label{eq:robinsecondcoefficient}
\end{align}
For a nonzero shift, the denominator is $\mathcal H_R(-\lambda)$ and the order is $\mu_\lambda$ from Equation~\eqref{eq:shiftedorder}. The extra order factor in Equation~\eqref{eq:robinprefactor} only changes $\kappa$ to $\kappa+1$ in the algebraic prefactor of Equation~\eqref{eq:shiftlogterm}; it does not change its logarithmic $1/\mu$ term. Consequently an identity first forces $x=4\pi$ and $\kappa=5/4$, then the leading normalization, and then $\lambda=0$ from the coefficient of $(\log\mu)/\mu$. This argument uses the shifted normalization even when the unshifted operator is not invertible.

We may therefore use Equation~\eqref{eq:robinlogratio}. Its $1/\mu$ coefficient forces
\begin{equation}
 h=h_*:=2\delta_{5/4}(4\pi)
       =2\pi^2-5\pi+\frac{11}{8}.
 \label{eq:matchedrobinh}
\end{equation}
Even granting $C_R=1$, the remaining $1/\mu^2$ coefficient is
\begin{align}
 B_*&:=b_{5/4,h_*}(4\pi)
     =\frac{13}{64}+\frac{5\pi}{4}-\pi^2+\frac{h_*^2}{8}\notag\\
 &=\frac{x^4-20x^3+90x^2-60x+225}{512}\bigg|_{x=4\pi}<0.
 \label{eq:robinanalyticobstruction}
\end{align}
For an explicit rational sign bound, $157/50<\pi<22/7$ and monotonicity of $2p^2-5p+11/8$ on that interval give
\begin{equation}
 B_*<\frac{13}{64}+\frac54\frac{22}{7}
           -\left(\frac{157}{50}\right)^2
           +\frac18\left(\frac{2123}{392}\right)^2
      =-\frac{1583907247}{768320000}<0.
 \label{eq:robinrationalbound}
\end{equation}
Thus the logarithm cannot vanish identically. No numerical Whittaker evaluation is a premise of the exclusion.
\end{proof}

The parameter tuple that survives the earlier filters is genuinely admissible, rather than excluded by a hidden positivity failure. Under the stated \emph{left-endpoint} convention $f'(u_0)=hf(u_0)$, the quadratic form is
\begin{equation}
 a_h[f]=\int_{u_0}^{\infty}\bigl(|f'|^2+V|f|^2\bigr)\dd u
                   +h|f(u_0)|^2.
 \label{eq:robinboundaryform}
\end{equation}
At $x=4\pi$, $\kappa=5/4$, and $h=h_*>0$, the potential satisfies
$V\geq16\pi(\pi-5/4)>0$ and the boundary term is nonnegative. Hence no adverse trace-inequality penalty is needed for this convention. Numerically $h_*\approx5.4062455342$ and $B_*\approx-2.0860522371$; Section~\ref{subsec:robindiagnostics} reports an independent normalization diagnostic. The failure is an exact second-order arithmetic mismatch after positivity, counting, scale, power, shift, and first-order matching have all survived.

Together, Corollary~\ref{cor:shiftedmorseexclusion} and Theorem~\ref{thm:robinmorseexclusion} cover every real constant separated self-adjoint left-endpoint condition for this Morse differential expression, including Neumann data as $h=0$. The right endpoint is limit point. The extension to a positive kinetic coefficient and a common exponential rate is proved in Section~\ref{subsec:kineticrate}; they do not leave a free spectral repair once the leading count is matched. Energy-dependent boundary data, additional potential terms, mode deletion, and a changed ratio of exponential rates remain outside this Morse-family exclusion.

\subsection{Kinetic--rate normalization of the full Morse family}
\label{subsec:kineticrate}
A positive kinetic coefficient and the common exponential rate do not leave an additional repair parameter when the ratio of the two rates remains $2:1$. Consider the full differential expression
\begin{equation}
 \mathcal A=-a\frac{\dd^2}{\dd u^2}
          +c^2\e^{2\beta u}-4c\kappa\e^{\beta u}+\lambda,
 \qquad u\geq u_0,\qquad a,c,\beta>0,
 \quad \kappa,u_0,\lambda\in\R.
 \label{eq:generalkineticmorse}
\end{equation}
The coordinate dilation and its unitary normalization are
\begin{equation}
 v=\frac{u}{\sqrt a},\qquad
 (\mathcal U_a f)(v)=a^{1/4}f(\sqrt a\,v),\qquad
 b:=\beta\sqrt a,\qquad x:=c\e^{\beta u_0}.
 \label{eq:kineticunitary}
\end{equation}
Here $\mathcal U_a:L^2([u_0,\infty),\dd u)\to
L^2([u_0/\sqrt a,\infty),\dd v)$ is unitary; no eigenvalue rescaling has been made in this first step.

\begin{corollary}[Full kinetic--rate Morse exclusion]
\label{cor:kineticratemorseexclusion}
No strictly positive realization of Equation~\eqref{eq:generalkineticmorse}, with a fixed real separated self-adjoint left-endpoint condition, has
$\det(\Id-w\mathcal A^{-1})=F(w)$. This includes arbitrary positive kinetic coefficient, common rate, potential scale, endpoint, real $\kappa$, constant spectral shift, and real constant Robin parameter.
\end{corollary}
\begin{proof}
Conjugation by $\mathcal U_a$ makes the kinetic coefficient one and replaces $\beta$ by $b$. A second change $r=bv+\log(2c/b)$, with unitary factor $b^{-1/2}$, gives
\begin{equation}
 \mathcal A\simeq b^2\left(-\frac{\dd^2}{\dd r^2}
             +\frac14\e^{2r}-\frac{2\kappa}{b}\e^r\right)+\lambda,
 \qquad r\geq r_0:=\log\frac{2x}{b}.
 \label{eq:generalmorsereduction}
\end{equation}
The first and second changes preserve the class of real constant separated boundary conditions. Apply the Morse counting law \cite[Theorem~3.2]{lagarias} with its parameter $k=-2\kappa/b$ and energy $(E-\lambda)/b^2$. For each fixed parameter tuple,
\begin{equation}
 N_{\mathcal A}(E)=\frac{\sqrt E}{b\pi}
       \left[\log\sqrt E+\log\frac2x-1\right]+O(1).
 \label{eq:kineticratecount}
\end{equation}
The fixed shift changes the displayed smooth terms only by $O(E^{-1/2}\log E)$, absorbed in this bounded remainder. The count also implies $\mathcal A^{-1}$ is trace class in the strictly positive case.

If its determinant were $F$, its complete eigenvalue count would satisfy Equation~\eqref{eq:requiredcount}. The two leading coefficients therefore force, successively,
\begin{equation}
 \frac{1}{2b\pi}=\frac{1}{4\pi},\qquad
 b=\beta\sqrt a=2,\qquad x=c\e^{\beta u_0}=4\pi.
 \label{eq:kineticratecalibration}
\end{equation}
When $b=2$, the first conjugated expression is exactly Equation~\eqref{eq:shiftedmorsedifferential} on $v\geq u_0/\sqrt a$. Dirichlet data stay Dirichlet, while $f'(u_0)=hf(u_0)$ becomes
\begin{equation}
 (\mathcal U_af)'(u_0/\sqrt a)
 =h_{\mathrm{red}}(\mathcal U_af)(u_0/\sqrt a),
 \qquad h_{\mathrm{red}}=\sqrt a\,h.
 \label{eq:kineticrobintransport}
\end{equation}
Corollary~\ref{cor:shiftedmorseexclusion} and Theorem~\ref{thm:robinmorseexclusion} now exclude every such normalized candidate. Unitary equivalence preserves the determinant, completing the proof.
\end{proof}

The first counting term fixes the product $\beta\sqrt a$, not $a$ and $\beta$ separately. The residual freedom with that product fixed is a coordinate dilation, not an untested spectral adjustment. In the Robin branch the last first-order match would read $\sqrt a\,h=h_*$; its second-order residual is still $B_*\neq0$. The remaining exclusions of scope are changes of the $2:1$ rate ratio, additional potential terms, mode deletion, and energy-dependent boundary data. These are different spectral families, not free parameters of Equation~\eqref{eq:generalkineticmorse}.

\section{Finite positivity does not close the infinite problem}
\label{sec:finite}
A finite matrix calculation can disprove a proposed positive representation if it produces a rigorously negative quadratic value. A positive calculation in finitely many degrees has a different logical status. The following construction makes the distinction explicit, including the possibility of off-line zeros arbitrarily far above the tested regime.

\begin{theorem}[Finite-tower nonclosure]
\label{thm:finite}
For every integer $q\geq0$ and every $0<\eta<1/2$, there exists a real entire function $G$ of order $1/2$ with $G(0)=1$ and $G(-x)>0$ for $x\geq0$, such that its logarithmic-moment matrices are positive definite for $0\leq N\leq q$, while $G$ has nonreal zeros. The associated even function $G(z^2)$ can have all its zeros in $|\Im z|<1/2$. The construction can also preserve the first two zeta counting terms.
\end{theorem}
\begin{proof}
Start with
\begin{equation}
 G_0(w)=\frac{\sin(\pi\sqrt w)}{\pi\sqrt w}
       =\prod_{n=1}^{\infty}(1-w/n^2).
 \label{eq:sinebase}
\end{equation}
Its logarithmic moments are $m_k^{(0)}=\zeta(2k+2)$, represented by the positive measure $\sum_{n\geq1}n^{-2}\delta_{n^{-2}}$. Every finite Hankel matrix is therefore positive definite by the infinite-support argument in Theorem~\ref{thm:moment}.

For $T>0$, put $a_T=(T+i\eta)^2$ and define
\begin{equation}
 G_T(w)=G_0(w)(1-w/a_T)(1-w/\overline{a_T}).
 \label{eq:finiteperturb}
\end{equation}
It has nonreal zeros $a_T,\overline{a_T}$, remains real entire of order $1/2$, and satisfies
$G_T(-x)=G_0(-x)|1+x/a_T|^2>0$. Its logarithmic moments are
\begin{equation}
 m_k^{(T)}=m_k^{(0)}+2\Re(a_T^{-k-1}).
 \label{eq:momentperturb}
\end{equation}
For $|a_T|\geq1$, the perturbation of the $(q+1)\times(q+1)$ Hankel matrix has operator norm at most $2(q+1)/|a_T|$. Choose $T$ large enough that this is less than the smallest eigenvalue of $H_q^{(0)}$. The perturbed $H_q$ remains positive definite, as do all its leading principal submatrices. The extra zeros of $G_T(z^2)$ are $\pm(T+i\eta)$ and $\pm(T-i\eta)$; all other zeros are real.

For the final assertion, replace $G_0$ by $D_{\mathrm{exp}}$. Its positive trace-class realization gives strictly positive finite moment matrices in exactly the same way. Multiplication by the same two factors adds only a fixed finite quartet in the $z$-plane and does not change either leading counting coefficient in Equation~\eqref{eq:expcount}.
\end{proof}

This theorem does not alter $F$ and does not construct an arithmetic counterexample. In particular, $G_T$ is not claimed to have the Euler product of $\zeta$. It shows why a proof of the specific arithmetic identity cannot be replaced by a finite list of matrix signs, strip symmetry, real-axis positivity, and two counting terms. Those properties coexist in the construction with nonreal zeros.

The theorem also does not exclude finite computation from a proof accompanied by an independent global estimate. A rigorously justified bound reducing the entire problem to a finite verification would be additional mathematical information. What fails is the unqualified inference from a finite positive prefix to positivity at every degree.

\subsection{Finite-to-infinite arithmetic moment transfer}
\label{subsec:momentlimit}
The finite-rank obstruction concerns one fixed representation, not a refinement family. A sequence can increase its rank while every individual member remains finite. The following consequence of Theorem~\ref{thm:moment} makes the required limiting statement explicit; it is a positivity-transfer result, not an independent proof of the arithmetic limit.

\begin{theorem}[Finite positive trace-moment transfer]
\label{thm:finitemomenttransfer}
RH is equivalent to the existence of positive finite-dimensional matrices $K_r$ such that, for every fixed $n\geq0$,
\begin{equation}
 \lim_{r\to\infty}\Tr K_r^{n+1}=m_n.
 \label{eq:finitemomentlimit}
\end{equation}
The matrices need not act on the same finite-dimensional space. No assumption of convergence uniform in $n$ is required.
\end{theorem}
\begin{proof}
For a fixed polynomial $p(t)=\sum_{j=0}^{N}c_jt^j$, positivity and finite-dimensional spectral calculus give
\begin{equation}
 \sum_{i,j=0}^{N}\overline{c_i}c_j\Tr K_r^{i+j+1}
 =\Tr\bigl(K_r p(K_r)^{*}p(K_r)\bigr)\geq0.
 \label{eq:finitetracegram}
\end{equation}
There are only finitely many terms for this $p$. Equation~\eqref{eq:finitemomentlimit} therefore gives $q_\xi(p)\geq0$. Since $p$ was arbitrary, Theorem~\ref{thm:moment} proves RH. Conversely, under RH the finite diagonal matrices with entries $\gamma_1^{-2},\ldots,\gamma_r^{-2}$ satisfy Equation~\eqref{eq:finitemomentlimit} by absolute convergence of Equation~\eqref{eq:powersums}. This conditional construction proves the equivalence; it is not an independent source construction.
\end{proof}

In Equation~\eqref{eq:finitetracegram}, self-adjointness gives $p(K_r)^*=\sum_i\overline{c_i}K_r^i$, which explains the conjugated coefficients on the left. The necessary information is the arithmetic value of every limiting moment. Conservation or refinement of an unrelated source operator cannot replace Equation~\eqref{eq:finitemomentlimit}. The result nevertheless shows why finite conservative discretizations remain legitimate ingredients of a proof. No single mesh is required to contain the complete arithmetic representation.

\begin{corollary}[Degree-wise error transfer]
\label{cor:degreeerror}
For a fixed $N$, suppose $K_r\geq0$ and
\begin{equation}
 \max_{0\leq n\leq2N}\left|\Tr K_r^{n+1}-m_n\right|
 \leq\epsilon_{r,N}.
 \label{eq:momenterror}
\end{equation}
Then $\lambda_{\min}(H_N)\geq-(N+1)\epsilon_{r,N}$. If $\epsilon_{r,N}\to0$ for every fixed $N$, RH follows.
\end{corollary}
\begin{proof}
The matrix $H_N^{(r)}=(\Tr K_r^{i+j+1})$ is positive semidefinite by Equation~\eqref{eq:finitetracegram}. The entrywise bound implies
$\|H_N-H_N^{(r)}\|\leq(N+1)\epsilon_{r,N}$. Apply this estimate to unit-vector quadratic forms and then take the limit. The final assertion is Theorem~\ref{thm:moment}.
\end{proof}

The bound in Equation~\eqref{eq:momenterror} must be proved against the \emph{actual} moments, not defined retrospectively from the candidate's computed values. At a finite $r$ it supplies an error statement, not all-degree positivity. No such all-degree arithmetic error estimate has been obtained for the source-compatible matrices in Equation~\eqref{eq:boundarycandidate}.

\subsection{A common Hilbert space and the determinant limit}
\label{subsec:tracenormlimit}
For an operator-limit route, specify isometric embeddings $\iota_r:\HS_r\to\HS$ into a common Hilbert space and write $\widetilde K_r=\iota_rK_r\iota_r^{*}$. The embeddings, not just the matrix entries, are part of the realization. Assume
\begin{equation}
 \|\widetilde K_r-K\|_1\longrightarrow0,
 \label{eq:tracenormlimit}
\end{equation}
where $\|\cdot\|_1$ is the trace norm. Completeness of the trace class and operator-norm convergence imply that $K$ is positive trace class. If $R$ bounds $\|\widetilde K_r\|$ and $\|K\|$, telescoping the powers gives
\begin{equation}
 \left|\Tr\widetilde K_r^{n+1}-\Tr K^{n+1}\right|
 \leq(n+1)R^n\|\widetilde K_r-K\|_1.
 \label{eq:powertraceerror}
\end{equation}
Indeed, write the difference of powers as the sum of $n+1$ products containing one factor $\widetilde K_r-K$ and use the trace-ideal inequality. Fredholm determinants also converge locally uniformly under Equation~\eqref{eq:tracenormlimit} \cite{simon,bornemann}. Neither assertion identifies the limit with $F$ unless the arithmetic moments or prime trace are independently matched.

Operator norm alone does not suffice. If $P_r$ projects onto the first $r$ coordinates of $\ell^2$ and $K_r=r^{-1}P_r$, then
\begin{equation}
 \|K_r\|\to0,\qquad \Tr K_r=1,
 \qquad \det(\Id-wK_r)=(1-w/r)^r\longrightarrow\e^{-w}.
 \label{eq:escapingtrace}
\end{equation}
The determinant of the operator-norm limit is one, not $\e^{-w}$. The example locates the failure in unresolved trace mass rather than in any individual large eigenvalue.

\subsection{A sufficient trace-tightness estimate}
\label{subsec:tracetightness}
A source approximation may first establish operator-norm convergence. The following lemma states an additional tail condition sufficient to upgrade it to the topology needed above. It uses only positive-operator block decomposition and Hilbert--Schmidt inequalities \cite{simon}.

\begin{lemma}[Uniform trace tightness]
\label{lem:tracetightness}
Suppose $T_r\geq0$ are trace-class operators on one Hilbert space, $\|T_r-T\|\to0$, and $\sup_r\Tr T_r\leq M<\infty$. Let $P_m$ be finite-rank orthogonal projections, strongly increasing to $\Id$, and put $Q_m=\Id-P_m$. If
\begin{equation}
 \sup_r\Tr(Q_mT_rQ_m)\leq\varepsilon_m,
 \qquad\varepsilon_m\downarrow0,
 \label{eq:tightnesstail}
\end{equation}
then $T$ is positive trace class and
\begin{equation}
 \|T_r-T\|_1
 \leq(\rank P_m)\|T_r-T\|
       +2\varepsilon_m+4\sqrt{M\varepsilon_m}.
 \label{eq:tightnessbound}
\end{equation}
In particular, $\|T_r-T\|_1\to0$.
\end{lemma}
\begin{proof}
Positivity passes to the operator-norm limit. For any finite-rank projection $P$, $\Tr(PTP)=\lim_r\Tr(PT_rP)\leq M$. Taking the supremum over such projections shows that $T$ is trace class with $\Tr T\leq M$. The same argument on the range of $Q_m$ gives $\Tr(Q_mTQ_m)\leq\varepsilon_m$.

For any positive trace-class $A$ with $\Tr A\leq M$ and $e=\Tr(Q_mAQ_m)$, the block decomposition gives
\begin{equation}
 \|A-P_mAP_m\|_1\leq e+2\sqrt{Me}.
 \label{eq:positiveblocktail}
\end{equation}
The diagonal tail has trace norm $e$. For the off-diagonal block, factor $P_mAQ_m=(P_mA^{1/2})(A^{1/2}Q_m)$ and bound its trace norm by the product of the Hilbert--Schmidt norms. Apply Equation~\eqref{eq:positiveblocktail} to both $T_r$ and $T$; bound the remaining finite block by $(\rank P_m)\|T_r-T\|$. This proves Equation~\eqref{eq:tightnessbound}. First choose $m$ to control the tails and then $r$ to control the finite block.
\end{proof}

The projections in Equation~\eqref{eq:tightnesstail} belong to the common Hilbert realization. A tail bound on a different, independently chosen basis at every refinement does not supply this hypothesis. The geometric preflights do not establish it, and a small conservation residual is not a substitute. The source-domain, excision, interface, and spectral limits must retain the declared order \cite{tcgsconservative,tcgsfoam}.

\subsection{Classical sign counts and delayed finite detection}
\label{subsec:chebotarev}
Chebotarev's complement gives a useful classical comparison to finite-prefix nonclosure. For the order-one real entire-function setting of Section~\ref{subsec:grommernormalization}, let
\begin{equation}
 \Delta_{-1}(f)=\Delta_0(f)=1,\qquad
 \Delta_n(f)=\det\left(\sum_\ell z_\ell^{-i-j-2}\right)_{i,j=0}^{n-1}
 \quad(n\geq1),
 \label{eq:chebotarevsequence}
\end{equation}
where the nonzero zeros $z_\ell$ are counted with multiplicity. The formulation in \cite[Theorem~10(ii)]{bariczstampach} states that, if the complete sequence
$\{\Delta_{n-1}(f)\Delta_n(f)\}_{n\geq0}$ contains exactly a finite number $m$ of negative terms, the function has $m$ \emph{distinct} nonreal conjugate pairs and infinitely many real zeros. This is a global sign-count statement, not an assertion that a finite positive prefix determines the total. We use it only for classical context in the nondegenerate sign-count setting; vanishing leading determinants require separate treatment and are not assigned a sign through rounding. No root-simplicity or determinant-nondegeneracy hypothesis is added to Theorem~\ref{thm:moment}.

For the construction of Theorem~\ref{thm:finite}, the order-one comparison function is $f_T(z)=G_T(z^2)$, not the order-$1/2$ function $G_T$ itself. The construction can delay failure in \emph{both} moment towers. For either stated positive base $G_0$, the shifted matrices are positive definite, since they are represented by the positive measure with weights $\lambda_j^2$ at its infinitely many distinct inverse zeros $\lambda_j>0$. Equation~\eqref{eq:momentperturb} gives, for $|a_T|\geq1$,
\begin{equation}
 \|H_q^{(T)}-H_q^{(0)}\|\leq\frac{2(q+1)}{|a_T|},\qquad
 \|H_q^{(+),(T)}-H_q^{(+),(0)}\|\leq\frac{2(q+1)}{|a_T|^2}.
 \label{eq:twotowerperturbation}
\end{equation}
Choosing $T$ sufficiently large makes both perturbations smaller than the corresponding least eigenvalues. The even/odd block identity then gives $\Delta_n(f_T)>0$ for every $n\leq2q+2$, although $f_T$ has the nonreal quartet $\pm(T+i\eta)$ and $\pm(T-i\eta)$. That quartet comprises two distinct conjugate pairs in the $z$-plane, whereas $G_T$ has one conjugate pair in the $w$-plane. Thus the first loss of strict positivity in the full lifted tower can be postponed past any prescribed finite prefix. Neither an eventual first negative index nor an exact global Chebotarev count is inferred from the truncated calculations; a zero determinant must also be distinguished from a negative one.

\section{Reproducible numerical diagnostics}
\label{sec:numerics}
\subsection{Computing without zero ordinates}
The coefficients used here are computed directly from derivatives at the central point:
\begin{equation}
 a_k=\frac{\xi^{(k)}(1/2)}{k!},\qquad
 f_j=(-1)^j\frac{a_{2j}}{a_0},
 \label{eq:taylorcoeff}
\end{equation}
followed by Equation~\eqref{eq:recurrence}. No table of zeta zeros is supplied to the calculation. This prevents the numerical input from presupposing the real spectral data that the programme is trying to derive.

Calculations used Python 3.13.5 and \texttt{mpmath} 1.3.0, with independent runs at 40 and 60 decimal working precision. The displayed digits agree between those runs. The software provides arbitrary-precision arithmetic, but these calculations use its ordinary floating-point context rather than interval enclosures \cite{mpmath}. The values are therefore reproducible numerical diagnostics, not certified sign proofs.

\begin{table}[htbp]
\centering
\caption{Logarithmic moments from central derivatives. Values are rounded, non-interval-certified diagnostics.}
\label{tab:moments}
\begin{tabular}{rl}
\toprule
$n$ & $m_n$\\
\midrule
0 & $2.31049931154190\times10^{-2}$\\
1 & $3.71725992852697\times10^{-5}$\\
2 & $1.44173931400973\times10^{-7}$\\
3 & $6.63031680252991\times10^{-10}$\\
4 & $3.21366415061660\times10^{-12}$\\
5 & $1.58769345149102\times10^{-14}$\\
6 & $7.90332572404674\times10^{-17}$\\
7 & $3.94647139999885\times10^{-19}$\\
8 & $1.97326507374497\times10^{-21}$\\
9 & $9.87216538329716\times10^{-24}$\\
\bottomrule
\end{tabular}
\end{table}

To improve the scale of the finite matrices, put $b_n=100^{n+1}m_n$ and $\widetilde H_N=(b_{i+j})$. If $D_N=\operatorname{diag}(1,100,\ldots,100^N)$, then
\begin{equation}
 \widetilde H_N=100D_NH_ND_N,\qquad
 \det\widetilde H_N=100^{(N+1)^2}\det H_N.
 \label{eq:scaling}
\end{equation}
This positive congruence preserves all sign questions. It is numerical conditioning, not an additional arithmetic or physical parameter.

\begin{table}[htbp]
\centering
\caption{Leading determinants after the exact positive scalings in Equations~\eqref{eq:scaling} and \eqref{eq:shiftedtowerscaling}. The prefixes are compatible with positivity; they do not certify either infinite tower. Both columns are evaluated through $N=4$.}
\label{tab:determinants}
\begin{tabular}{rll}
\toprule
$N$ & $\det\widetilde H_N$ & $\det\widetilde H_N^{(+)}$\\
\midrule
0 & $2.31049931154190$ & $3.71725992852697\times10^{-1}$\\
1 & $1.94933555481914\times10^{-1}$ & $3.86048846787094\times10^{-3}$\\
2 & $2.17281051052625\times10^{-4}$ & $3.11996028999615\times10^{-7}$\\
3 & $1.17758171601017\times10^{-9}$ & $9.24213895905459\times10^{-14}$\\
4 & $1.96218205320998\times10^{-17}$ & $7.86624627750459\times10^{-23}$\\
\bottomrule
\end{tabular}
\end{table}

The rapid decrease of the determinants makes precision control important. Agreement between two precision settings is a practical diagnostic of stability, not a substitute for an enclosure of truncation and roundoff errors. A formal numerical certificate would need validated derivatives or validated theta integrals, interval propagation through Equation~\eqref{eq:recurrence}, and certified matrix positivity, for example through an interval $LDL^*$ factorization. No claim of such certification is made here.

\subsection{Independent checks of the failed candidate}
The exponential asymptotic can also be inspected directly. Let $c_{\mathrm{exp}}$ be the positive constant in Equation~\eqref{eq:expmismatch}. The quotient
\begin{equation}
 R_{\mathrm{exp}}(s)=
 \frac{F(-(s-1/2)^2)}{c_{\mathrm{exp}}s^{9/4}
 D_{\mathrm{exp}}(-(s-1/2)^2)}
\end{equation}
has numerical values approximately $2.6975168723$, $1.2678702215$, and $1.0597685216$ at $s=20,80,320$, respectively. Their approach toward one is a diagnostic of the derived asymptotic, not its proof. The proof is Equation~\eqref{eq:besselbound} and the gamma-ratio calculation.

For the Morse repair, $\delta_{9/4}\approx-1.8300625401$. The quantity
$\mu\bigl(F/D_{9/4}/C_{9/4}-1\bigr)$, evaluated at $s=2\mu+1/2$, is approximately $-1.63031121$, $-1.76862626$, and $-1.80943692$ for $\mu=30,100,300$. These values agree with the limiting coefficient in Equation~\eqref{eq:morseratio}. Again, the family-level exclusion depends on the analytic nonzero coefficient, not on the floating-point values.

\subsection{Quantitative prime-trace comparison}
\label{subsec:traceerrors}
A numerical comparison becomes an exclusion test only after its errors are bounded. Write, for real $s>1$ and $t=s-1/2$,
\begin{align}
 \mathcal T_K(s)&=2t\Tr\bigl[K(\Id+t^2K)^{-1}\bigr],\label{eq:Tdefinition}\\
 B_\xi(s)&=\frac1s+\frac1{s-1}-\frac12\log\pi
                 +\frac12\psi_\Gamma(s/2),\notag\\
 \mathcal A_\xi(s)&=B_\xi(s)-\sum_{n=2}^{\infty}\Lambda(n)n^{-s}.
 \label{eq:Adefinition}
\end{align}
The target is $\mathcal T_K=\mathcal A_\xi$, not the prime sum without its completion factors. This comparison uses the convergent half-plane. The critical strip is $0<\Re s<1$ and the critical line is $\Re s=1/2$; analytic continuation is well defined there, while the original absolutely convergent Euler product is not. The choice $s>1$ controls this series manipulation, not a supposed instability of complex analysis \cite{titchmarsh,dlmf}.

\begin{proposition}[Positive resolvent-trace bound]
\label{prop:resolventbound}
For positive trace-class $K,Q$ on the same Hilbert space and real $s>1$,
\begin{equation}
 |\mathcal T_K(s)-\mathcal T_Q(s)|
 \leq2(s-\tfrac12)\|K-Q\|_1.
 \label{eq:resolventtracebound}
\end{equation}
No commutation between $K$ and $Q$ is required.
\end{proposition}
\begin{proof}
Put $a=t^2$. Multiplication on the left by $\Id+aK$ and on the right by $\Id+aQ$ verifies
\begin{equation}
 K(\Id+aK)^{-1}-Q(\Id+aQ)^{-1}
 =(\Id+aK)^{-1}(K-Q)(\Id+aQ)^{-1}.
 \label{eq:noncommutingresolvent}
\end{equation}
Both resolvents have operator norm at most one by positivity. The trace-ideal inequality bounds the trace norm of the right side by $\|K-Q\|_1$. Apply $|\Tr R|\leq\|R\|_1$ and multiply by $2t$.
\end{proof}

\begin{lemma}[Unconditional prime-power tail]
\label{lem:primetail}
For an integer $P\geq3$ and real $s>1$,
\begin{equation}
 0\leq\sum_{n>P}\frac{\Lambda(n)}{n^s}
 \leq \mathcal R_{\mathrm{pr}}(P,s)
 :=P^{1-s}\left(\frac{\log P}{s-1}+\frac1{(s-1)^2}\right).
 \label{eq:primetail}
\end{equation}
\end{lemma}
\begin{proof}
The defining prime-power weights satisfy $0\leq\Lambda(n)\leq\log n$. The function $x\mapsto(\log x)x^{-s}$ is decreasing for $x\geq3$. Bound the sum by $\int_P^\infty(\log x)x^{-s}\dd x$ and integrate. No zero-location assumption enters.
\end{proof}

Suppose a positive trace-class limit $K$ has a declared positive approximation $\widetilde K_r$ with a proved bound $\|\widetilde K_r-K\|_1\leq\eta_r$. Define the exact finite discrepancy
\begin{equation}
 \Delta_{r,P}(s)=\mathcal T_{\widetilde K_r}(s)
 -\left[B_\xi(s)-\sum_{n=2}^{P}\Lambda(n)n^{-s}\right].
 \label{eq:finitediscrepancy}
\end{equation}
Let $\widehat\Delta_{r,P}(s)$ be its computed value with certified absolute evaluation error at most $\epsilon_{\mathrm{eval}}$. Matrix assembly, quadrature, and source regularization must either be included in $\eta_r$ or bounded separately in the evaluation error; no contribution is silently set to zero.

\begin{corollary}[Certified rejection and limiting closure]
\label{cor:certifiedtrace}
If, at one real $s>1$,
\begin{equation}
 |\widehat\Delta_{r,P}(s)|>
 2(s-\tfrac12)\eta_r+\mathcal R_{\mathrm{pr}}(P,s)
 +\epsilon_{\mathrm{eval}},
 \label{eq:tracerejection}
\end{equation}
then $K$ does not satisfy Equation~\eqref{eq:primebridge}. Conversely, fix $1<s_-<s_+$ and suppose $\eta_r\to0$, $P_r\to\infty$, the evaluation error tends uniformly to zero, and
\begin{equation}
 \sup_{s\in[s_-,s_+]}|\widehat\Delta_{r,P_r}(s)|\longrightarrow0.
 \label{eq:uniformtraceclosure}
\end{equation}
Then $K$ satisfies Equation~\eqref{eq:primebridge} on the interior of that interval and RH follows.
\end{corollary}
\begin{proof}
Proposition~\ref{prop:resolventbound} and Lemma~\ref{lem:primetail} give
\begin{equation}
 \left|\mathcal T_K(s)-\mathcal A_\xi(s)
                   -\widehat\Delta_{r,P}(s)\right|
 \leq2(s-\tfrac12)\eta_r+\mathcal R_{\mathrm{pr}}(P,s)
       +\epsilon_{\mathrm{eval}}.
 \label{eq:totaltraceerror}
\end{equation}
The reverse triangle inequality proves rejection. On the compact interval, the first term is at most $2(s_+-1/2)\eta_r$ and the prime tail is at most $\mathcal R_{\mathrm{pr}}(P,s_-)$. The limiting assumptions therefore imply equality throughout the interval. Apply Theorem~\ref{thm:primetrace}.
\end{proof}

A finite grid of apparently small discrepancies does not establish Equation~\eqref{eq:uniformtraceclosure}; it requires an independently proved bound between grid points. Likewise, a computable prime tail does not provide the geometric trace-norm error. The latter remains a substantive spectral convergence problem for the candidate source construction.

\subsection{Exclusion of a bounded operator family}
\label{subsec:operatorfamily}
The certified-cover requirement of the conservative paper can be specialized to arithmetic discrepancy \cite{tcgsconservative}. Let $\Theta$ be a compact metric parameter set, with a certified finite $\delta$-net $\{\theta_j\}_{j=1}^{J}$, and suppose positive trace-class candidates satisfy
\begin{equation}
 \|K_\theta-K_{\theta'}\|_1\leq L\,d(\theta,\theta')
 \label{eq:familylipschitz}
\end{equation}
for a proved uniform constant $L$. Fix one $s>1$. Let $\widehat\Delta_j$ approximate $\mathcal T_{K_{\theta_j}}(s)-\mathcal A_\xi(s)$ with certified total error $E_j$, obtained for example from Equation~\eqref{eq:totaltraceerror}. Then
\begin{equation}
 \min_j\bigl(|\widehat\Delta_j|-E_j\bigr)
       -2(s-\tfrac12)L\delta>0
 \label{eq:familyrejection}
\end{equation}
excludes every $K_\theta$ in the declared family. Indeed, choose a net point within distance $\delta$ of $\theta$ and use Equation~\eqref{eq:resolventtracebound} to bound the change in discrepancy. Neither compactness alone nor a finite sample of parameters supplies $L$ or a certified cover. Changing the family or its boundary law after rejection defines a new construction rather than repairing the rejected identity.

\subsection{Exact manufactured algebraic diagnostics}
\label{subsec:exactdiagnostics}
A small rational example verifies the elimination, normalization, and trace algebra without claiming to realize the source geometry. In the interior--boundary order, take
\begin{equation}
 \mathsf C_*=
 \begin{pmatrix}2&-1&-1\\-1&2&0\\-1&0&2\end{pmatrix},
 \quad M_{b,*}=\Id,
 \quad\mathsf S_*=
 \begin{pmatrix}3/2&-1/2\\-1/2&3/2\end{pmatrix},
 \quad K_*=\begin{pmatrix}3/4&1/4\\1/4&3/4\end{pmatrix}.
 \label{eq:rationalexample}
\end{equation}
The full matrix has positive leading minors $2,3,4$. The Schur response has eigenvalues $1,2$, and $K_*$ has eigenvalues $1,1/2$. Consequently
\begin{equation}
 \Tr K_*^{n+1}=1+2^{-n-1},\qquad
 \det H^{(*)}_1=\frac1{8},\quad\det H^{(*)}_2=0.
 \label{eq:rationalmoments}
\end{equation}
The rank-two termination is an exact control for Theorem~\ref{thm:rank}, not a value of an arithmetic Hankel determinant. The source and static boundary determinants are visibly different:
\begin{equation}
 \frac{\det(\mathsf C_*-w\Id)}{\det\mathsf C_*}
 =\frac{(2-w)((2-w)^2-2)}4,
 \quad
 \det(\Id-wK_*)=(1-w)(1-w/2).
 \label{eq:rationaldeterminants}
\end{equation}
This directly verifies the need to distinguish them in Equation~\eqref{eq:bulkboundarydeterminant}.

The accompanying \texttt{boundary\_trace\_diagnostics.py} uses exact rational arithmetic to check the Schur complement, minimizing extension, inverse normalization, finite moments, rank termination, and spectral-pencil factorization, together with the noncommuting resolvent identity at several rational parameters. Appendix~\ref{app:boundarycode} gives a minimal implementation. These checks verify algebra and code only. They are not a four-dimensional solver, a constructed singular-boundary functional, a trace-norm convergence certificate, or evidence for RH. The original numerical moment and Morse diagnostics remain unchanged.


\subsection{Central normalization and the quarter-shift check}
\label{subsec:quartershift}
The familiar zero sum at $s=1$ is related to, but not equal to, the central logarithmic moment $m_0$. The Laurent expansion of $\zeta$ at one and the gamma values in Equation~\eqref{eq:xi} give \cite{titchmarsh,dlmf}
\begin{equation}
 \frac{\xi'(1)}{\xi(1)}
 =1+\frac{\gamma_{\mathrm E}}2-\frac12\log(4\pi)
 =M(-\tfrac14)
 =\frac12\sum_\rho\frac1{\rho(1-\rho)}.
 \label{eq:quartershift}
\end{equation}
Indeed, $\rho(1-\rho)=w_\rho+1/4$ and differentiation of
$F(-(s-1/2)^2)=\xi(s)/\xi(1/2)$ gives
$\xi'(s)/\xi(s)=2(s-1/2)M(-(s-1/2)^2)$. The sum in Equation~\eqref{eq:quartershift} includes all nontrivial zeros with multiplicity and converges absolutely. In contrast,
\begin{equation}
 m_0=M(0)=\frac{\xi''(1/2)}{2\xi(1/2)},\qquad
 m_0-M(-\tfrac14)
 =\frac14\sum_j\frac1{w_j(w_j+1/4)}.
 \label{eq:quartershiftdifference}
\end{equation}
These identities do not assume RH. Numerically,
$M(-1/4)\approx0.0230957089661210$ and
$m_0-M(-1/4)\approx9.28414929794\times10^{-6}$, consistently with Table~\ref{tab:moments}. The distinction is a shift of spectral argument, not a discrepancy between two evaluations of the same moment. The added \texttt{scale\_endpoint\_diagnostics.py} reproduces these values and checks the free-endpoint comparison coefficients; its numerical values, like the earlier special-function diagnostics, are not interval certificates.


\subsection{Robin normalization and second-order diagnostics}
\label{subsec:robindiagnostics}
At the parameter tuple forced by the counting, power, shift, and first inverse-order conditions, the Whittaker normalization can be evaluated directly. The following evaluation uses $\kappa=5/4$, $x=4\pi$, $\lambda=0$, and $h=h_*$ from Equation~\eqref{eq:matchedrobinh}. The derivative is computed both by differentiation of $W$ and from
$W_{\kappa,\nu}(x)=\e^{-x/2}x^{\nu+1/2}U(\nu+1/2-\kappa,1+2\nu,x)$ with $\partial_xU(a,b,x)=-aU(a+1,b+1,x)$ \cite[Sections~13.3 and 13.14]{dlmf}.

\begin{equation}
\begin{aligned}
 h_*&\approx 5.4062455342297510,\\
 W_{5/4,0}(4\pi)&\approx 0.0422104279006116,\\
 \partial_xW_{5/4,0}(4\pi)&\approx-0.0167498199338303,\\
 Q_{5/4,h_*,0}(4\pi)&\approx-0.6913792552608059,\\
 C_R&\approx1.0446407439095589,\\
 B_*&\approx-2.0860522370547263.
\end{aligned}
\label{eq:robinvalues}
\end{equation}
The displayed special-function values are non-interval-certified diagnostics; the exclusion uses the exact rational sign bound in Equation~\eqref{eq:robinrationalbound}, not the numerical value of $C_R$.

The normalization therefore also differs numerically from one, but no claim of a certified special-function enclosure is required. More decisively, the numerical values of
\begin{equation}
 \mu^2\log\!\left(\frac{F(-4\mu^2)}{C_RD_R(-4\mu^2)}\right)
 \label{eq:robinsecondnumerical}
\end{equation}
are approximately $-1.8759731882$, $-2.0296884422$, $-2.0679420328$, and $-2.0806912775$ at $\mu=30,100,300,1000$, respectively. They approach the analytically derived $B_*$ with the $O(\mu^{-1})$ remainder implied by Equation~\eqref{eq:robinlogratio}. The script \texttt{robin\_shift\_diagnostics.py} repeats the calculation at declared working precision, compares the two derivative formulas, and verifies the coefficient polynomials and rational sign bound by exact arithmetic. The original numerical outputs remain separate and unchanged.

\subsection{The shifted Stieltjes prefix and block diagnostics}
\label{subsec:shiftedprefix}
The shifted column in Table~\ref{tab:determinants} makes the Hamburger--Stieltjes comparison numerically explicit without changing the original unshifted results. With $b_n=100^{n+1}m_n$ and the diagonal $D_N$ from Equation~\eqref{eq:scaling}, the corresponding normalization is
\begin{equation}
 \widetilde H_N^{(+)}=(b_{i+j+1})_{i,j=0}^{N}
   =100^2D_NH_N^{(+)}D_N,\qquad
 \det\widetilde H_N^{(+)}
   =100^{(N+1)(N+2)}\det H_N^{(+)}.
 \label{eq:shiftedtowerscaling}
\end{equation}
The additional factor of $100$ per matrix entry relative to the unshifted congruence is essential. Both transformations preserve the exact sign conditions.

The script \texttt{shifted\_tower\_diagnostics.py} recomputes the central derivatives and recurrence at 40 and 60 decimal working digits. It does not read a zero table or reuse rounded moments as numerical inputs. The four shifted determinants, for $N=0,1,2,3$, agree at every digit printed in Table~\ref{tab:determinants}. The same script compares the directly assembled scaled Grommer determinant with its even/odd block product for sizes $1$ through $8$. The companion script \texttt{completed\_tower\_diagnostics.py} uses central Taylor order $4N+4=20$ at $N=4$ to compute $m_9$ and complete both columns through degree four. Fresh runs at 40 and 60 decimal working digits agree on every displayed moment and determinant digit; the direct even/odd block checks cover sizes $1$ through $10$. These are floating-point consistency tests, not interval certificates or evidence for all-degree positivity.

The separate standard-library script \texttt{kinetic\_rate\_diagnostics.py} performs exact-rational checks of the reduction coefficients in Equation~\eqref{eq:generalmorsereduction}, the Robin transport, and the scaling and block identities on manufactured finite data. These implementation checks are distinguished from the symbolic all-parameter proof in Corollary~\ref{cor:kineticratemorseexclusion} and from the unresolved arithmetic realization.

\section{Invariants, artifacts, assumptions, and identifications}
\label{sec:classification}
The framework's invariant language must be indexed by an actual comparison class. The following distinctions prevent a numerical convention or an attractive spectrum from acquiring an unsupported source interpretation \cite{tcgsfoundation,tcgsoperator}.

\begin{longtable}{P{0.17\textwidth}P{0.35\textwidth}P{0.40\textwidth}}
\caption{Logical status of the principal structures.}\label{tab:types}\\
\toprule
Class & Structure & Meaning and boundary\\
\midrule
\endfirsthead
\toprule
Class & Structure & Meaning and boundary\\
\midrule
\endhead
Governing source architecture & $\CS^4$, $\mathcal S=\Orb(p_0)$, typed ECL, foliation--label--clock distinction & The starting architecture of the TCGS application; not inferred from the signs in Tables~\ref{tab:moments}--\ref{tab:determinants}.\\[0.5em]
Arithmetic invariance & Zeros and multiplicities of the fixed entire function $F$; exact moments $m_n$ & Exact mathematical data. A representation preserving them is not thereby identical to the source ontology.\\[0.5em]
Selector invariance & Unitary transport in Equation~\eqref{eq:naturality} & Preserves the Gram form along declared equivalence arrows; not across arbitrary changes of potential or endpoint.\\[0.5em]
Coordinate convention & $w=z^2$, inverse spectral variable, half-density normalization & Must be retained consistently; changing variables is not a change in the arithmetic target.\\[0.5em]
Numerical artifact & Working precision, rounding, scale factor $100$, finite matrix size & Influences computation, not the exact sign condition. The factor $100$ is an explicit congruence.\\[0.5em]
Model assumption & Exponential potential, Dirichlet endpoint, Morse correction $\kappa$ & Independently testable choices. Their calibration to a counting law is not an ECL derivation.\\[0.5em]
Unproved identification & $q_\xi(p)=\|L_\sigma J_\sigma p\|^2$ for every polynomial & The sufficient all-degree bridge. Its existence has not been obtained from the source geometry.\\[0.5em]
Excluded identification & A zeta zero is automatically a member of $\mathcal S$ & A zero of an entire function is not a singularity of that function. Its logarithmic derivative has a pole, which is still a different mathematical type from a source singularity.\\[0.5em]
Proved model exclusion & The complete inverse determinant is not $F$ for any strictly positive real constant separated Morse realization in Equation~\eqref{eq:generalkineticmorse}. & Free scale/endpoint: Corollary~\ref{cor:scaleendpointexclusion}; constant shift: Corollary~\ref{cor:shiftedmorseexclusion}; Robin: Theorem~\ref{thm:robinmorseexclusion}; kinetic/rate: Corollary~\ref{cor:kineticratemorseexclusion}. This is a family-specific exclusion, not a general spectral no-go theorem.\\
\bottomrule
\end{longtable}

Source invariants and their mathematical representatives must also remain distinct. The prime-power coefficients $\Lambda(n)$ are not declared to be direct measurements of Counterspace. They enter as exact arithmetic data on the right side of the bridge identity. If a source representation is eventually obtained, the identification will be a theorem relating those data to the representation, not a renaming of primes as singularities.

The same care applies to positivity. Positive definiteness of $H_N$ is a mathematical property of a finite matrix. Positive semidefiniteness of all $H_N$ is an infinite arithmetic assertion equivalent to RH. Positivity of $O_\sigma$ is a theorem about a closed readout operator. These statements are related only through a proved identity such as Equation~\eqref{eq:grambridge}; their shared adjective does not establish that identity.

\section{Discussion}
\label{sec:discussion}
\subsection{What the geometric attempt establishes}
The analysis moves the TCGS application from a broad source--shadow analogy to a sharply typed realization problem. The arithmetic target is fixed independently through $\xi$. The positive geometry is fixed independently through operator domains, norms, and quadratic forms. The Gram and prime-trace theorems specify the exact equality required between them. The construction is therefore not allowed to succeed merely by displaying the critical half, a real spectrum, or a visually plausible distribution of levels.

The explicit corridor model demonstrates both the opportunity and the obstruction. Noncompact geometry supplies the logarithmic increase in accessible length needed for the zeta density. This explains concretely why a logarithmic counting term need not require a temporal source dimension or an ordinary compact spectrum. Yet the fully computed determinant fails. A first repair within the Morse family removes the leading power defect only at one parameter and then fails at the next coefficient. A constant Robin endpoint supplies an additional power and one further matching coefficient, so it requires a separate calculation: even after that matching and exclusion of constant shifts, Equation~\eqref{eq:robinanalyticobstruction} leaves a nonzero second-order term. The failure is not attributed to insufficient numerical precision; it is a proved mismatch of analytic functions.

The bounded exclusion is useful for subsequent source-geometry work. A candidate in the full Morse family of Equation~\eqref{eq:generalkineticmorse}, with any real constant separated endpoint condition, cannot satisfy Equation~\eqref{eq:primebridge}, however its parameters are interpreted. A different candidate must change the spectral realization in a mathematically substantive way. Possible changes include boundary data beyond a real constant separated condition, a nonlocal constitutive reduction, singular interface data, or a different operator class. These are design directions, not derived solutions, and each must reproduce the exact arithmetic identity rather than only the two counting coefficients.

\subsection{The distinction between a criterion and a solution}
An equivalence theorem does not establish either side of the equivalence unconditionally. Theorem~\ref{thm:moment} says that universal positivity settles RH; it does not prove universal positivity. Theorem~\ref{thm:primetrace} says that a positive realization of the prime trace settles RH; it does not produce that realization. Their value is to identify a more explicit mathematical target, not to relabel the original assertion as a completed proof.

This distinction is particularly important for a framework with a complete source ontology. Completeness of source content is not the same proposition as completeness of a chosen mathematical invariant or faithfulness of a selected representation. The framework itself distinguishes the full profile from one readout and an observability operator from a Hamiltonian \cite{tcgsfoundation,tcgsoperator}. The present analysis applies those distinctions to arithmetic rather than weakening them at the point where a desired theorem becomes available.

The mathematical conclusions also do not depend on assigning physical ontology to the auxiliary gamma measures, Mellin coordinate, or heat parameter. Those objects are devices in analysis. The source framework constrains how a physical realization would be interpreted, while the proofs remain statements about explicitly defined functions and operators. This division avoids both an ontological inflation of mathematics and an exemption of framework-inspired mathematics from ordinary proof obligations.

\subsection{A decisive next construction}
The most direct constructive target is a positive trace-class $K$ specified from an admissible source reduction, together with a derivation of Equation~\eqref{eq:primebridge} in $s>1$. This target has two advantages. It does not require prior knowledge of the zero ordinates, and its prime-power series is absolutely convergent. A successful derivation could use an independently defined arithmetic action, trace formula, or constitutive insertion map. None has been supplied by the constructions in this paper.

A second target is an all-degree proof that the logarithmic transform of the exact theta moments satisfies Equation~\eqref{eq:grambridge}. The raw theta kernel is known explicitly, and the coefficient recurrence is executable. The missing result is a positive representation of the nonlinear logarithmic moment transform. Establishing positivity of the raw theta density again would not address this remaining problem.

The boundary conditions for success are exact. The operator must be positive on its stated domain; the trace must be finite where used; multiplicities and the full determinant must agree; selector changes claimed to be equivalent must preserve the construction; and no unknown arithmetic fact may be inserted as a spectral assumption. A finite positive prefix, an asymptotic density match, or an operator built from already selected real zeros does not meet these conditions.

\subsection{Limitations of the present results}
No constructive arithmetic realization of $\Info$, $\Read_{\mathcal S,G}$, or $J_\sigma$ has been established here. No positivity-preservation theorem for the full logarithmic moment transform has been proved. The numerical matrices are finite and non-interval-certified. The compact-elliptic obstruction applies only to the operator class stated in Proposition~\ref{prop:weylobstruction}, and the Morse exclusion to the full constant-coefficient differential expression in Equation~\eqref{eq:generalkineticmorse}, with rate ratio $2:1$, arbitrary real constant spectral shift, and fixed real separated left-endpoint condition, as proved in Corollary~\ref{cor:kineticratemorseexclusion} using Corollary~\ref{cor:shiftedmorseexclusion} and Theorem~\ref{thm:robinmorseexclusion}. Neither is a general impossibility theorem for spectral approaches.

The foundational and application papers on gravity, biology, quantum readout, and consciousness provide the source architecture and its type discipline. Their domain-specific empirical claims are not used as premises for an arithmetic theorem. The three constituent framework works cited here are those needed for source geometry, typed Counterduction, and operator realization. This restriction keeps the mathematical dependence explicit without treating unrelated successes or analogies as evidence for RH.

Finally, the moment formulation is grounded in the classical Grommer--Hamburger real-zero method, with its reciprocal-square normalization distinguished explicitly in Section~\ref{subsec:grommernormalization} \cite{grommer,schmudgen,bariczstampach}. The Bessel--Morse spectral family has established antecedents \cite{lagarias}. The paper's results should be evaluated as an integrated constructive investigation with explicit transfer conditions and exclusions, not as a claim to have invented moment theory, Fredholm determinants, or spectral interpretations of zeta functions.

\subsection{What the conservative boundary construction adds}
\label{subsec:discussionboundary}
The conservative and reciprocal-foam constructions change the available geometric starting point without supplying the missing arithmetic identity. Compatible cochains give an explicit positive frozen reciprocal sector. A shared boundary functional, once derived, would place the singular channels in the same variational assembly. Positive interior elimination then produces a boundary response and supplies a naturality theorem under the stated selector transports. These are constructive statements at a declared mathematical level; they do not turn $\mathsf S_r$ or $K_r$ into the ontic ECL \cite{tcgsconservative,tcgsfoam}.

Two distinctions now prevent a misleading success claim. First, the static boundary determinant need not equal the bulk determinant: Equation~\eqref{eq:bulkboundarydeterminant} displays the missing interior factor and spectral dependence. Second, a smooth compact boundary response remains subject to a Weyl obstruction, even though it is nonlocal. The proposed use of interfaces is therefore a method of generating and testing candidates, not a declaration that a boundary construction automatically possesses the arithmetic density.

The finite-to-infinite result addresses a different gap. Theorem~\ref{thm:finitemomenttransfer} shows that finite matrices can participate in a proof when their traces converge to every actual arithmetic moment. Lemma~\ref{lem:tracetightness} gives sufficient spectral-tail control for a trace-class operator limit, while Equation~\eqref{eq:totaltraceerror} separates geometric approximation, prime truncation, and numerical evaluation. The resulting exclusions are deterministic mathematical comparisons; they do not require a physical laboratory experiment. A finite positive prefix still cannot replace the stated limits.

The strongest remaining construction is an independently specified source representation whose boundary or bulk operator satisfies either Equation~\eqref{eq:finitemomentlimit} or Equation~\eqref{eq:primebridge}. The geometric papers do not derive this arithmetic representation, their shared singular-boundary functional, or the trace-tightness bounds required here. Their liquid-lens, curvature, magnetic, and horizon sectors retain their own purposes and limits. They are not premises of the number-theoretic proof and are not imported as arithmetic evidence \cite{tcgsconservative,tcgsfoam}.

The arithmetic content is specific. The prime-power term can be written
\begin{equation}
 \sum_{n=2}^{\infty}\Lambda(n)n^{-s}
 =\sum_{p\ \mathrm{prime}}\sum_{k=1}^{\infty}
          (\log p)\e^{-sk\log p},\qquad s>1.
 \label{eq:primitiveprimedata}
\end{equation}
A proposed corridor interpretation must account for these particular weights, repetitions, and lengths, as well as the completion terms in $B_\xi(s)$. General recurrence, self-similarity, or nontrivial topology does not do so. Arithmetic inputs may legitimately be declared in a mathematical model; declaring them is not the same as deriving their source representation. The restriction is against hiding the desired zero-location conclusion in the operator definition.

\section{Conclusion}
\label{sec:conclusion}
The moment and spectral formulations of the Riemann hypothesis become directly comparable when their arithmetic and positivity requirements are separated and then connected by an exact identity. In the normalization used here,
\begin{equation}
 \mathrm{RH}
 \ \Longleftrightarrow\ H_N\succeq0\quad\text{for all }N
 \ \Longleftrightarrow\ F(w)=\det(\Id-wK),\quad K\geq0\text{ trace class}.
 \label{eq:conclusionsummary}
\end{equation}
The prime-trace identity in Theorem~\ref{thm:primetrace} expresses the same target entirely within the absolutely convergent prime-power half-plane. The constitutive Gram-transfer theorem specifies what a TCGS arithmetic realization must preserve, and the finite-rank obstruction shows that no single finite protected sector can supply the complete representation.

The geometric construction has been carried far enough to be rejected precisely. A positive exponential corridor reproduces the leading arithmetic density but has the wrong determinant. A one-parameter Morse correction removes its power mismatch only at $\kappa=9/4$, where a nonzero inverse-order coefficient still forbids equality. Corollary~\ref{cor:scaleendpointexclusion} extends the exclusion to free potential scale and endpoint: their only spectral combination is fixed to $4\pi$ by the second counting coefficient, so neither supplies a remaining repair parameter. These are mathematical results of the attempted construction, not evidence that all geometric routes fail.

The Morse-family exclusion now also permits arbitrary potential scale, endpoint, constant spectral shift, and real constant separated endpoint condition. Corollary~\ref{cor:kineticratemorseexclusion} additionally permits positive kinetic coefficient and common exponential rate: the first counting term fixes their product $\beta\sqrt a=2$, and the second fixes $c\e^{\beta u_0}=4\pi$. After the scale and power are fixed, the shift is forced to zero; the Robin parameter can remove the first-order term but cannot remove the explicit second-order residual. This conclusion uses neither a numerical estimate of the normalization constant nor an extrapolation from a finite list of eigenvalues.

The unresolved step is now explicit: derive, without assuming RH, a source-grounded positive representation of the full logarithmic moment form or the exact prime-weighted resolvent trace. No such derivation is established in this manuscript, and the Riemann hypothesis is therefore not claimed as solved. The result is a rigorous arithmetic realization programme with a proved criterion, an executed model family, and exact conditions separating a genuine proof from a spectral analogy.

The compatible-source extension supplies a further conditional route: positive interior elimination yields finite boundary candidates, coefficientwise convergence of their traces to the arithmetic moments implies RH, and trace tightness controls an operator limit when that route is used. The static/bulk determinant distinction and the regular boundary-power obstruction prevent artificial spectral closure. The rational checks establish the finite algebra exactly, whereas the shared source boundary law and the all-degree arithmetic identification remain unconstructed.

\appendix
\section{Derivation of the Whittaker comparison coefficient}
\label{app:whittaker}
This appendix proves Lemma~\ref{lem:whittaker}, including the term used to exclude the only power-matched Morse candidate. The integral representation \cite{dlmf}, for fixed $x>0$ and sufficiently large real $\mu$, is
\begin{equation}
 W_{\kappa,\mu}(x)=
 \frac{x^{\mu+1/2}2^{-2\mu}}{\Gamma(\mu+1/2-\kappa)}
 \int_1^\infty\e^{-xt/2}
 (t-1)^{\mu-1/2-\kappa}(t+1)^{\mu-1/2+\kappa}\dd t.
 \label{eq:whittakerintegral}
\end{equation}
After $y=xt/2$, this becomes
\begin{equation}
 W_{\kappa,\mu}(x)=
 \frac{x^{1/2-\mu}\Gamma(2\mu)}{\Gamma(\mu+1/2-\kappa)}
 \mathbb E\!\left[R_\mu(Y)\mathbf1_{\{Y>x/2\}}\right],
 \label{eq:gammaexpectation}
\end{equation}
where $Y$ has gamma shape $2\mu$, unit scale, and
\begin{equation}
 R_\mu(y)=
 \left(1-\frac{x^2}{4y^2}\right)^{\mu-1/2}
 \left(\frac{1+x/(2y)}{1-x/(2y)}\right)^\kappa.
 \label{eq:Rmu}
\end{equation}
The expectation is an auxiliary integral notation; it makes no claim about physical randomness.

On $y\geq\mu$ and for $\mu$ sufficiently large, Taylor expansion gives uniformly
\begin{equation}
 R_\mu(y)=1+\frac{\kappa x}{y}
 -\frac{(\mu-1/2)x^2}{4y^2}+O(\mu^{-2}).
 \label{eq:Rexpansion}
\end{equation}
Indeed, the logarithm has the two displayed terms with remainder $O(\mu/y^4+y^{-3})$, while the square of its leading part is $O(\mu^{-2})$ on this region.

The omitted region contributes an exponentially small error. To see this without an uncontrolled expansion near $y=x/2$, put $v=x/(2y)\in(0,1)$. If $\mu-1/2\geq|\kappa|$, the two factors $(1-v)^{\mu-1/2-\kappa}(1+v)^{\mu-1/2+\kappa}$ imply $R_\mu(y)\leq2^{2|\kappa|}$. Moreover, for $Y$ of shape $2\mu$,
\[
 \Pr(Y\leq\mu)\leq\e^{\mu}\mathbb E\e^{-Y}
 =\left(\frac{\e}{4}\right)^\mu.
\]
The analogous negative-moment tails are exponentially small as well, by applying the same bound to gamma shapes $2\mu-1$ and $2\mu-2$ and multiplying by their gamma-ratio factors. Thus the full moments may replace the truncated moments in Equation~\eqref{eq:Rexpansion}.

Using $\mathbb E Y^{-1}=(2\mu-1)^{-1}$ and
$\mathbb E Y^{-2}=((2\mu-1)(2\mu-2))^{-1}$, we obtain
\begin{equation}
 \mathbb E\!\left[R_\mu(Y)\mathbf1_{\{Y>x/2\}}\right]
 =1+\frac{\kappa x/2-x^2/16}{\mu}+O(\mu^{-2}).
 \label{eq:expectationexpansion}
\end{equation}
The duplication formula and the ordinary gamma-ratio expansion give
\begin{align}
 \frac{x^{1/2-\mu}\Gamma(2\mu)}{\Gamma(\mu+1/2-\kappa)}
 &=\frac{\sqrt x}{2\sqrt\pi}\Gamma(\mu)(x/4)^{-\mu}
   \frac{\Gamma(\mu+1/2)}{\Gamma(\mu+1/2-\kappa)}\notag\\
 &=\frac{\sqrt x}{2\sqrt\pi}\Gamma(\mu)(x/4)^{-\mu}\mu^\kappa
   \left[1-\frac{\kappa^2}{2\mu}+O(\mu^{-2})\right].
 \label{eq:gammaratioappendix}
\end{align}
Multiplication by Equation~\eqref{eq:expectationexpansion} proves Equations~\eqref{eq:whittakerexpansion}--\eqref{eq:whittakercoefficient}. In particular, the coefficient excluding $\kappa=9/4$ is not inferred from a numerical fit to Whittaker values.


\subsection{Second-order coefficients and uniform differentiation}
\label{subsec:secondorderwhittaker}
The next coefficient can be extracted from the same integral without fitting special-function values. Put $\alpha=\kappa x$, $\beta=x^2/4$, and let $Y$ have the auxiliary gamma distribution of shape $2\mu$ from Equation~\eqref{eq:gammaexpectation}. Uniformly on $y\geq\mu$ and on each fixed compact positive $x$-interval,
\begin{align}
 R_\mu(y)
 &=1+\frac\alpha y-\frac{(\mu-1/2)\beta}{y^2}
   +\frac12\left(\frac\alpha y-\frac{\mu\beta}{y^2}\right)^2
   +O(\mu^{-3}).
 \label{eq:secondRexpansion}
\end{align}
This follows by expanding $\log(1-x^2/(4y^2))$ and
$\log((1+x/(2y))/(1-x/(2y)))$ and then exponentiating. The omitted logarithmic terms are $O(\mu y^{-4}+y^{-3})$, and the cubic exponential terms are $O(\mu^{-3})$ on this region. The same estimates hold after one $x$-derivative.

For the complementary region $x/2<y<\mu$, the earlier gamma tail bound remains valid. The derivative also has an integrable uniform estimate: with $v=x/(2y)$,
\begin{equation}
 \partial_xR_\mu(y)
 =\frac{R_\mu(y)}{1-v^2}
       \left(\frac\kappa y-\frac{(\mu-1/2)x}{2y^2}\right).
 \label{eq:Rderivativebound}
\end{equation}
For $\mu-3/2\geq|\kappa|$, the first factor is bounded by $2^{2|\kappa|}$ by the same endpoint-factor argument used in Appendix~\ref{app:whittaker}. The remaining factor is $O(1+\mu)$ uniformly for $x$ in a compact positive interval. The resulting tail and derivative tail are exponentially small. The moving lower endpoint contributes no term for sufficiently large $\mu$, since $R_\mu(x/2)=0$. Negative-moment tails are controlled by the corresponding shifted gamma shapes.

Using
$\mathbb E Y^{-j}=\prod_{\ell=1}^{j}(2\mu-\ell)^{-1}$ for $j=1,2,3,4$, Equation~\eqref{eq:secondRexpansion} gives
\begin{align}
 e_1&=\frac\alpha2-\frac\beta4,\notag\\
 \mathbb E\!\left[R_\mu(Y)\mathbf1_{\{Y>x/2\}}\right]
 &=1+\frac{e_1}{\mu}
     +\frac{(\alpha-\beta)/4+e_1^2/2}{\mu^2}
     +O(\mu^{-3}),\notag\\
 \log\mathbb E\!\left[R_\mu(Y)\mathbf1_{\{Y>x/2\}}\right]
 &=\frac{\kappa x/2-x^2/16}{\mu}
     +\frac{\kappa x/4-x^2/16}{\mu^2}
     +O(\mu^{-3}).
 \label{eq:secondgammaexpectation}
\end{align}
The bounds just proved justify differentiation of this asymptotic with respect to $x$, rather than differentiation of an uncontrolled pointwise remainder.

For fixed $a,b$, the logarithmic gamma-ratio expansion is
\begin{equation}
 \log\frac{\Gamma(\mu+a)}{\Gamma(\mu+b)}
 =(a-b)\log\mu
  +\frac{B_2(a)-B_2(b)}{2\mu}
  -\frac{B_3(a)-B_3(b)}{6\mu^2}+O(\mu^{-3}),
 \label{eq:bernoulligammaratio}
\end{equation}
where $B_2(t)=t^2-t+1/6$ and $B_3(t)=t^3-3t^2/2+t/2$ \cite[Section~5.11]{dlmf}. With $a=1/2$, $b=1/2-\kappa$, this becomes
\begin{equation}
 \log\frac{\Gamma(\mu+1/2)}{\Gamma(\mu+1/2-\kappa)\mu^\kappa}
 =-\frac{\kappa^2}{2\mu}
   +\frac{-\kappa^3/6+\kappa/24}{\mu^2}+O(\mu^{-3}).
 \label{eq:secondwhittakergamma}
\end{equation}
Combining with Equation~\eqref{eq:secondgammaexpectation} proves Equation~\eqref{eq:whittakerlogsecond}. Its differentiated version gives Equation~\eqref{eq:whittakerlogderivative}. Independently, substituting $a=1/4$, $b=0$ in Equation~\eqref{eq:bernoulligammaratio} gives $-3/(32\mu)-1/(128\mu^2)$. Adding
$\log(1-1/(16\mu^2))=-1/(16\mu^2)+O(\mu^{-4})$ from the elementary factor in $\xi(2\mu+1/2)$ gives $-9/(128\mu^2)$, proving Equation~\eqref{eq:xilogsecond}; the zeta correction is exponentially small.

Finally, taking the logarithm of Equation~\eqref{eq:robinprefactor} contributes
$h/(2\mu)+(-\kappa x/2+x^2/8-h^2/8)/\mu^2$ to the Robin characteristic. Subtraction from the target yields the coefficient $b_{\kappa,h}(x)$ in Equation~\eqref{eq:robinsecondcoefficient}. The nonzero residual in Equation~\eqref{eq:robinanalyticobstruction} is therefore obtained from the same controlled integral and gamma expansions as the Dirichlet coefficient, not from a numerical estimate of $C_R$.

\section{A minimal reproducibility implementation}
\label{app:code}
The following self-contained code reproduces the unshifted moment and determinant prefixes through $m_8$. Its Taylor order is $4N+2$, because $H_N$ requires $m_0,\ldots,m_{2N}$ and the recurrence requires $f_1,\ldots,f_{2N+1}$. It contains no calls to routines that return zeta zeros. The separate source package includes a command-line implementation, precision-specific JSON results, and the Morse asymptotic diagnostic.

\begin{lstlisting}[language=Python,caption={Finite logarithmic-moment diagnostics; not an interval certificate.}]
import mpmath as mp

# Run once at 40 and once at 60 decimal working precision.
mp.mp.dps = 60
N = 4
scale = mp.mpf(100)

def xi(s):
    return (s * (s - 1) / 2 * mp.pi**(-s/2)
            * mp.gamma(s/2) * mp.zeta(s))

a = mp.taylor(xi, mp.mpf('0.5'), 4*N + 2)
f = [(-1)**j * a[2*j] / a[0] for j in range(2*N + 2)]
m = []
for n in range(2*N + 1):
    value = -(n + 1) * f[n + 1]
    value -= mp.fsum(f[k] * m[n-k] for k in range(1, n+1))
    m.append(value)

b = [scale**(n+1) * value for n, value in enumerate(m)]
for n, value in enumerate(m):
    print('moment', n, mp.nstr(value, 25))
for degree in range(N + 1):
    H = mp.matrix([[b[i+j] for j in range(degree+1)]
                   for i in range(degree+1)])
    print('determinant', degree, mp.nstr(mp.det(H), 25))

def morse_check(mu):
    k = mp.mpf(9) / 4
    x = 4 * mp.pi
    s = 2 * mu + mp.mpf('0.5')
    central = xi(mp.mpf('0.5'))
    C = 2 * mp.whitw(k, 0, x) / (central * mp.pi**mp.mpf('.25'))
    target = xi(s) / central
    candidate = mp.whitw(k, mu, x) / mp.whitw(k, 0, x)
    return mu * (target / candidate / C - 1)

for mu in (30, 100, 300):
    print('Morse correction', mu, mp.nstr(morse_check(mp.mpf(mu)), 20))
print('Analytic limit', mp.nstr(mp.pi**2 - 9*mp.pi/2
                               + mp.mpf(39)/16, 20))
\end{lstlisting}

For rigorous numerical certification, each displayed scalar would need an enclosure rather than a precision-dependent approximation. Such an upgrade would strengthen the finite diagnostics only. Theorem~\ref{thm:finite} would still prevent the certified finite prefix from becoming an all-degree result without a separate global argument.

\clearpage
\section{Notation and proof-dependence summary}
\label{app:notation}
\begin{longtable}{P{0.25\textwidth}P{0.68\textwidth}}
\caption{Symbols whose formal types must remain distinct.}\\
\toprule
Symbol & Meaning\\
\midrule
\endfirsthead
\toprule
Symbol & Meaning\\
\midrule
\endhead
$\CS^4$, $\Sigma^3$ & Counterspace source representation and physical Shadow carrier.\\
$\src$, $M_{\src}$ & Admissible source configuration and its informational organization.\\
$\ES$, $\mathcal E_{\mathcal S,\src}$ & Constitutive map and complete profile value.\\
$\sigma$, $\mathcal O^\sigma_{\mathcal S}$ & Readout selector and physical-readout map.\\
$L_\sigma$, $O_\sigma$ & Closed readout derivative and observability operator $L_\sigma^*L_\sigma$.\\
$\xi(s)$, $\Xi(z)$ & Riemann completion and its critical-line-centered entire function.\\
$F(w)$ & Normalized even-power reduction; no square-root branch is chosen.\\
$w_j$, $\lambda_j$ & Zeros of $F$; under RH, inverse squared ordinates $\lambda_j=\gamma_j^{-2}$.\\
$A_{2j}$, $f_j$ & Raw theta moments and Taylor coefficients of $F$.\\
$m_n$, $H_N$, $q_\xi$ & Logarithmic moments, finite Hankel matrices, and their polynomial form.\\
$J_\sigma$ & Unconstructed arithmetic insertion map in the Gram-transfer condition.\\
$K$, $A$ & Positive trace-class arithmetic candidate and its inverse on the nonzero spectral subspace.\\
$A_{\mathrm{exp}}$, $A_\kappa$ & Explicit exponential and calibrated Dirichlet Morse test operators.\\
$D_{\mathrm{exp}}$, $D_\kappa$ & Their normalized Fredholm determinants; neither equals $F$.\\
$\mathcal A$, $a$, $\beta$ & General Morse realization in Equation~\eqref{eq:generalkineticmorse}, its positive kinetic coefficient, and the lower exponential rate.\\
$b$, $x$, $\lambda$ & Kinetic--rate invariant $b=\beta\sqrt a$, scale--endpoint combination $x=c\e^{\beta u_0}$, and constant spectral shift.\\
$h$, $h_{\mathrm{red}}$ & Original left-endpoint Robin coefficient and its unitary transport $h_{\mathrm{red}}=\sqrt a\,h$.\\
$\mathcal H_D$, $\mathcal H_R$ & Unshifted Dirichlet and Robin characteristic functions in Equation~\eqref{eq:robincharacteristic}; not yet normalized determinants.\\
$Q_{\kappa,h,\nu}(x)$ & Robin characteristic combination $2x\partial_xW_{\kappa,\nu}(x)-(1+h)W_{\kappa,\nu}(x)$.\\
$D_R$, $C_R$ & Normalized Robin determinant at $\lambda=0$ and its asymptotic normalization constant; shifted normalization uses Equation~\eqref{eq:shiftedcharacteristicproduct}.\\
$h_*$, $B_*$ & The first-order matched Robin coefficient and the nonzero second-order obstruction in Equations~\eqref{eq:matchedrobinh} and \eqref{eq:robinanalyticobstruction}.\\
$H_N^{(+)}$, $\widetilde H_N^{(+)}$ & Shifted Stieltjes moment matrix and its positively scaled numerical version.\\
$\Delta_n(f)$ & Classical inverse-power Hankel determinant for an order-one function $f$; it is not the diagonal numerical scaling $D_N$.\\
$\Lambda(n)$, $\psi_\Gamma$ & Von Mangoldt prime-power weights and gamma logarithmic derivative.\\
$\tau$, $\mu$, $\kappa$ & Heat probe parameter, large Whittaker order, and Morse model parameter; none is an ontic temporal dimension.\\
\bottomrule
\end{longtable}

\begin{samepage}
The proof dependence can be read without the framework interpretation. No unconditional analytic or operator result in this paper uses the TCGS axioms as a proof premise. The framework supplies the realization setting and admissibility restrictions; the conditional transfer statements require the separately displayed mathematical hypotheses. Standard analytic properties of $\xi$ give Lemma~\ref{lem:product}; the Hamburger theorem then gives Theorem~\ref{thm:moment}. Positive trace-class theory gives Theorems~\ref{thm:determinant} and \ref{thm:primetrace}. A \emph{proved} Gram identification would invoke these through Theorem~\ref{thm:gramtransfer}; that identification remains unconstructed. Independently, classical spectral theory and the explicit calculations give Propositions~\ref{prop:weylobstruction} and \ref{prop:heat}, the corridor construction, and Theorems~\ref{thm:expmismatch} and \ref{thm:morseexclusion}. Corollary~\ref{cor:scaleendpointexclusion} uses the same Whittaker expansion and the scale--endpoint counting law. Lemma~\ref{lem:complexwhittakergrowth} and Proposition~\ref{prop:whittakerproducts} justify the complex-plane product step; Corollary~\ref{cor:shiftedmorseexclusion} and Theorem~\ref{thm:robinmorseexclusion} use the displayed shift and second-order asymptotics, independently of numerical normalization values; Equations~\eqref{eq:grommerparity}--\eqref{eq:grommerdeterminants} relate the arithmetic indexing to the classical criterion. Corollary~\ref{cor:kineticratemorseexclusion} uses a unitary dilation and the two leading counting terms to reduce the full Morse family to the already excluded one. Section~\ref{subsec:chebotarev} gives classical sign-count context and the two-tower finite-prefix comparison; the shifted numerical column uses Equation~\eqref{eq:shiftedtowerscaling}. The finite numerical results are not premises in any of these proofs.\par
\end{samepage}

\section{Exact boundary-response verification}
\label{app:boundarycode}
The following independent rational check reproduces the inverse, moment sequence, and first rank-termination conditions in Equations~\eqref{eq:rationalexample}--\eqref{eq:rationalmoments}, and verifies the noncommuting resolvent identity. It needs no external package. The complete command-line script additionally checks the minimizing extension, full characteristic polynomial, and spectral-pencil factorization, and writes \texttt{boundary\_trace\_diagnostics.json}. All 28 checks in that script pass using exact fractions. These manufactured checks do not evaluate the source boundary functional or the Riemann prime trace.

\begin{lstlisting}[language=Python,caption={Exact rational verification of the manufactured boundary example.}]
from fractions import Fraction as F

I = [[F(1), F(0)], [F(0), F(1)]]
S = [[F(3, 2), F(-1, 2)], [F(-1, 2), F(3, 2)]]
K = [[F(3, 4), F(1, 4)], [F(1, 4), F(3, 4)]]
Q = [[F(3, 4), F(0)], [F(0), F(1, 4)]]

def mul(A, B):
    return [[sum((A[i][k]*B[k][j] for k in range(2)), F(0))
             for j in range(2)] for i in range(2)]

def add(A, B, c=F(1)):
    return [[A[i][j] + c*B[i][j] for j in range(2)]
            for i in range(2)]

def det(A):
    return A[0][0]*A[1][1] - A[0][1]*A[1][0]

def inv(A):
    d = det(A)
    if not d:
        raise ValueError('Singular matrix')
    return [[A[1][1]/d, -A[0][1]/d],
            [-A[1][0]/d, A[0][0]/d]]

def check(condition):
    if not condition:
        raise ArithmeticError('Exact identity failed')

check(mul(S, K) == I)
check(mul(K, Q) != mul(Q, K))
P, moments = I, []
for n in range(5):
    P = mul(P, K)
    moments.append(P[0][0] + P[1][1])
    check(moments[-1] == 1 + F(1, 2**(n+1)))
a, b, c, d, e = moments
check(a*c - b*b == F(1, 8))
check(a*(c*e-d*d)-b*(b*e-c*d)+c*(b*d-c*c) == 0)
for t in (F(1), F(3, 2), F(4)):
    RK, RQ = inv(add(I, K, t*t)), inv(add(I, Q, t*t))
    left = add(mul(K, RK), mul(Q, RQ), F(-1))
    right = mul(mul(RK, add(K, Q, F(-1))), RQ)
    check(left == right)
print('Exact manufactured checks passed; no RH claim.')
\end{lstlisting}

\section*{Data, code, and computational assistance}
No experimental dataset is used. The numerical inputs are the defining special functions, exact constants, and parameter choices printed in the manuscript. The source package contains \texttt{main.tex}, \texttt{diagnostics.py}, \texttt{morse\_diagnostics.py}, and machine-readable outputs. It also contains \texttt{robin\_shift\_diagnostics.py}, which checks the added coefficient algebra and reproduces the Robin normalization and large-order comparisons. 

The package additionally contains \path{scale_endpoint_diagnostics.py}, \path{shifted_tower_diagnostics.py}, \path{completed_tower_diagnostics.py}, \path{kinetic_rate_diagnostics.py}, and \path{sesquilinear_diagnostics.py}, with their machine-readable outputs. The driver \path{verify_release.py} reruns all nine diagnostic scripts in a fresh directory and checks the regenerated JSON files against the packaged outputs and the printed values in Tables~\ref{tab:moments}--\ref{tab:determinants} and Equations~\eqref{eq:robinvalues}--\eqref{eq:robinsecondnumerical}.

Language-model assistance was used in developing the mathematical argument, drafting, preparing code, and checking LaTeX consistency. This disclosure is not a claim of independent peer verification; responsibility for final verification and submission rests with the named author.

The geometric source papers are cited as manuscripts at their stated fourth-amendment dates; no repository identifier is invented for them. The additional script \texttt{boundary\_trace\_diagnostics.py} and its JSON output use exact rational arithmetic and are separate from the unchanged floating-point xi and Morse diagnostics. The amended source package also includes a preservation report and a source manifest distinguishing source-derived statements, proved conditional implications, and unexecuted arithmetic closure.

\section*{License}
Copyright \copyright\ 2026 Henry Arellano-Pe\~na. This manuscript and its original source code are distributed under the Creative Commons Attribution 4.0 International License (CC BY 4.0): \url{https://creativecommons.org/licenses/by/4.0/}. Third-party works cited in the bibliography retain their own licenses.

\begin{thebibliography}{99}

\bibitem{riemann}
B. Riemann,
\emph{Ueber die Anzahl der Primzahlen unter einer gegebenen Gr\"osse},
Monatsberichte der Berliner Akademie (1859), 671--680.
Original text and translations: \url{https://www.maths.tcd.ie/pub/HistMath/People/Riemann/Zeta/}.

\bibitem{titchmarsh}
E. C. Titchmarsh,
\emph{The Theory of the Riemann Zeta-Function},
2nd ed., revised by D. R. Heath-Brown,
Clarendon Press, Oxford, 1986.

\bibitem{clay}
Clay Mathematics Institute,
\emph{Riemann Hypothesis}, Millennium Prize Problems,
\url{https://www.claymath.org/millennium/riemann-hypothesis/},
accessed 5 September 2026.

\bibitem{tcgsfoundation}
H. Arellano-Pe\~na,
\emph{Timeless Counterspace \& Shadow Gravity---A Unified Framework:
4D-Counterduction, Foundational Consistency, and Cartographic Inquiry},
Preprints, version 2, 11 August 2026, not peer reviewed.
\url{https://doi.org/10.20944/preprints202511.1472.v2}.

\bibitem{tcgsgeometry}
H. Arellano,
\emph{Source--Shadow Geometry of TCGS--SEQUENTION:
Gravito-Capillary Foams, a Retrocausal Non-Local Foliation,
and a Disciplined Cartographic Program for Cluster Offsets,
Cosmic Anisotropy, and Biological Invariants},
v3.0 integrated, closure-audited preprint, June 2026.
\url{https://www.researchgate.net/publication/397427092}.

\bibitem{tcgsoperator}
H. Arellano-Pe\~na,
\emph{Discrete Records Alone Do Not Determine Gauge Ontology:
Connes Rigidity, Lattice Continuum Limits, and Source--Shadow
Reconstruction in TCGS--SEQUENTION},
preprint, 5 August 2026.
\url{https://doi.org/10.13140/RG.2.2.33830.89921}.

\bibitem{berrykeating}
M. V. Berry and J. P. Keating,
The Riemann zeros and eigenvalue asymptotics,
\emph{SIAM Review} \textbf{41} (1999), 236--266.
\url{https://doi.org/10.1137/S0036144598347497}.

\bibitem{li}
X.-J. Li,
The positivity of a sequence of numbers and the Riemann hypothesis,
\emph{Journal of Number Theory} \textbf{65} (1997), 325--333.

\bibitem{cardon}
D. A. Cardon and S. A. Roberts,
An equivalence for the Riemann Hypothesis in terms of orthogonal polynomials,
\emph{Journal of Approximation Theory} \textbf{138} (2006), 54--64.
\url{https://doi.org/10.1016/j.jat.2005.09.016}.

\bibitem{voros}
A. Voros,
\emph{Zeta Functions over Zeros of Zeta Functions},
Lecture Notes of the Unione Matematica Italiana, vol. 8,
Springer, 2010.
\url{https://doi.org/10.1007/978-3-642-05203-3}.

\bibitem{grommer}
J. Grommer,
Ganze transzendente Funktionen mit lauter reellen Nullstellen,
\emph{Journal f\"ur die reine und angewandte Mathematik} \textbf{144} (1914), 114--166.
\url{https://doi.org/10.1515/crll.1914.144.114}.
Archival record: \url{https://eudml.org/doc/149420}.

\bibitem{schmudgen}
K. Schm\"udgen,
\emph{The Moment Problem},
Graduate Texts in Mathematics, vol. 277,
Springer, Cham, 2017.
\url{https://doi.org/10.1007/978-3-319-64546-9}.

\bibitem{simon}
B. Simon,
\emph{Trace Ideals and Their Applications},
2nd ed., Mathematical Surveys and Monographs, vol. 120,
American Mathematical Society, Providence, 2005.

\bibitem{reed}
M. Reed and B. Simon,
\emph{Methods of Modern Mathematical Physics, I: Functional Analysis},
revised and enlarged ed., Academic Press, New York, 1980.

\bibitem{hormander}
L. H\"ormander,
The spectral function of an elliptic operator,
\emph{Acta Mathematica} \textbf{121} (1968), 193--218.
\url{https://doi.org/10.1007/BF02391913}.

\bibitem{lagarias}
J. C. Lagarias,
The Schr\"odinger operator with Morse potential on the right half-line,
\emph{Communications in Number Theory and Physics} \textbf{3} (2009), 323--361.
\url{https://doi.org/10.4310/CNTP.2009.v3.n2.a3}.
Preprint: \url{https://arxiv.org/abs/0712.3238}.

\bibitem{dlmf}
National Institute of Standards and Technology,
\emph{Digital Library of Mathematical Functions},
Chapters 5, 10, 13, and 25,
\url{https://dlmf.nist.gov/}, accessed 5 September 2026.

\bibitem{bombieri}
E. Bombieri,
\emph{Problems of the Millennium: The Riemann Hypothesis},
Clay Mathematics Institute, official problem description.
\url{https://www.claymath.org/wp-content/uploads/2022/05/riemann.pdf}.

\bibitem{mpmath}
The mpmath development team,
\emph{mpmath: A Python Library for Arbitrary-Precision Floating-Point Arithmetic},
software version 1.3.0 used in the computations.
\url{https://mpmath.org/}.


\bibitem{tcgsconservative}
H. Arellano-Pe\~na,
\emph{Conservative Discrete Source-to-Readout Maps in Four Dimensions:
Accountability Conditions, Median-Dual Geometry, Singular-Orbit Topology,
and Boundary Completion, with a TCGS--SEQUENTION Application},
manuscript, fourth amendment, 28 August 2026, CC BY 4.0.

\bibitem{tcgsfoam}
H. Arellano-Pe\~na,
\emph{Reciprocal Fractal-Foam Geometry of Counterspace:
Core Geometry, Conservative Discretization, Domain-Merger Benchmark,
and Falsification Guide},
manuscript, fourth amendment, 28 August 2026, CC BY 4.0.

\bibitem{arnoldfeec}
D. N. Arnold, R. S. Falk, and R. Winther,
Finite element exterior calculus: from Hodge theory to numerical stability,
\emph{Bulletin of the American Mathematical Society} \textbf{47} (2010), 281--354.
\url{https://doi.org/10.1090/S0273-0979-10-01278-4}.
Preprint: \url{https://arxiv.org/abs/0906.4325}.

\bibitem{kron}
F. D\"orfler and F. Bullo,
\emph{Kron Reduction of Graphs with Applications to Electrical Networks},
preprint, 2011.
\url{https://arxiv.org/abs/1102.2950}.

\bibitem{steklov}
A. Girouard and I. Polterovich,
\emph{Spectral Geometry of the Steklov Problem},
preprint, 2014.
\url{https://arxiv.org/abs/1411.6567}.

\bibitem{bornemann}
F. Bornemann,
On the numerical evaluation of Fredholm determinants,
\emph{Mathematics of Computation} \textbf{79} (2010), 871--915.
\url{https://doi.org/10.1090/S0025-5718-09-02280-7}.
Preprint: \url{https://arxiv.org/abs/0804.2543}.

\bibitem{levin}
B. Ja. Levin,
\emph{Distribution of Zeros of Entire Functions},
revised ed., Translations of Mathematical Monographs, vol. 5,
American Mathematical Society, Providence, 1980.

\bibitem{bariczstampach}
\'A. Baricz and F. \v{S}tampach,
The Hurwitz-type theorem for the regular Coulomb wave function via Hankel determinants,
\emph{Linear Algebra and its Applications} \textbf{548} (2018), 259--272.
\url{https://doi.org/10.1016/j.laa.2018.03.012}.
Preprint: \url{https://arxiv.org/abs/1708.07729v3}.

\bibitem{cravencsordas}
T. Craven and G. Csordas,
Iterated Laguerre and Tur\'an inequalities,
\emph{Journal of Inequalities in Pure and Applied Mathematics}
\textbf{3} (2002), no. 3, Article 39, 14 pp.
\url{https://eudml.org/doc/122436}.
\end{thebibliography}
\end{document}
